\section{问题一逐用例结果}
\begin{footnotesize}
\setlength{\tabcolsep}{2.6pt}
\begin{longtable}{l r rr rr rr rr}
\caption{问题一逐用例结果：Makespan（周期）与总额外搬运量（字节）}\label{tab:app-p1}\\
\toprule
\multirow{2}{*}{用例} & \multirow{2}{*}{单核基准} & \multicolumn{2}{c}{2 核} & \multicolumn{2}{c}{3 核} & \multicolumn{2}{c}{4 核} & \multicolumn{2}{c}{5 核} \\
\cmidrule(lr){3-4}\cmidrule(lr){5-6}\cmidrule(lr){7-8}\cmidrule(lr){9-10}
 & & Makespan & 额外搬运 & Makespan & 额外搬运 & Makespan & 额外搬运 & Makespan & 额外搬运 \\
\midrule
\endfirsthead
\toprule
用例 & 单核基准 & \multicolumn{2}{c}{2 核} & \multicolumn{2}{c}{3 核} & \multicolumn{2}{c}{4 核} & \multicolumn{2}{c}{5 核} \\
\midrule
\endhead
\bottomrule
\endfoot
case\_001 & 233{,}110 & 116{,}868 & 1{,}152 & 93{,}135 & 862{,}464 & 58{,}984 & 3{,}456 & 47{,}502 & 4{,}608 \\
case\_002 & 261{,}945 & 133{,}944 & 5{,}165{,}568 & 101{,}425 & 964{,}608 & 77{,}554 & 766{,}464 & 65{,}085 & 910{,}848 \\
case\_003 & 602{,}212 & 372{,}446 & 6{,}590{,}892 & 280{,}827 & 7{,}330{,}768 & 267{,}227 & 8{,}010{,}228 & 256{,}255 & 8{,}971{,}928 \\
case\_004 & 392{,}357 & 213{,}333 & 4{,}306{,}944 & 134{,}043 & 221{,}184 & 108{,}126 & 3{,}741{,}696 & 95{,}289 & 3{,}852{,}288 \\
case\_005 & 95{,}618 & 74{,}654 & 2{,}341{,}728 & 67{,}164 & 2{,}443{,}372 & 67{,}076 & 2{,}693{,}956 & 67{,}076 & 2{,}693{,}956 \\
case\_006 & 82{,}837 & 39{,}720 & 39{,}190 & 30{,}961 & 94{,}344 & 30{,}961 & 94{,}344 & 30{,}961 & 94{,}344 \\
case\_007 & 1{,}012{,}562 & 579{,}668 & 11{,}993{,}088 & 338{,}094 & 55{,}296 & 338{,}094 & 55{,}296 & 203{,}834 & 110{,}592 \\
case\_008 & 487{,}605 & 244{,}266 & 0 & 163{,}359 & 0 & 123{,}060 & 0 & 100{,}603 & 0 \\
case\_009 & 153{,}261 & 92{,}198 & 172{,}032 & 92{,}198 & 172{,}032 & 92{,}198 & 172{,}032 & 92{,}198 & 172{,}032 \\
case\_010 & 86{,}969 & 45{,}272 & 48{,}384 & 33{,}937 & 836{,}616 & 29{,}586 & 362{,}880 & 29{,}586 & 362{,}880 \\
case\_011 & 184{,}697 & 92{,}646 & 180{,}224 & 69{,}625 & 1{,}081{,}344 & 46{,}880 & 540{,}672 & 46{,}880 & 540{,}672 \\
case\_012 & 59{,}088 & 30{,}094 & 28{,}800 & 21{,}894 & 20{,}160 & 17{,}253 & 39{,}168 & 14{,}300 & 49{,}536 \\
case\_013 & 263{,}949 & 132{,}198 & 8{,}192 & 88{,}351 & 16{,}384 & 66{,}982 & 24{,}576 & 54{,}363 & 32{,}768 \\
case\_014 & 17{,}698{,}626 & 8{,}813{,}761 & 81{,}618{,}360 & 5{,}934{,}584 & 81{,}312{,}936 & 4{,}512{,}560 & 94{,}384{,}536 & 3{,}677{,}086 & 86{,}530{,}056 \\
case\_015 & 177{,}362 & 92{,}411 & 224{,}536 & 60{,}766 & 278{,}280 & 54{,}989 & 528{,}396 & 41{,}506 & 317{,}952 \\
case\_016 & 7{,}715{,}523 & 7{,}704{,}912 & 393{,}264 & 6{,}868{,}180 & 119{,}561{,}444 & 5{,}445{,}426 & 119{,}559{,}028 & 5{,}445{,}426 & 119{,}559{,}028 \\
case\_017 & 162{,}810 & 75{,}539 & 92{,}160 & 59{,}088 & 179{,}712 & 39{,}520 & 263{,}808 & 34{,}720 & 347{,}904 \\
case\_018 & 1{,}004{,}367 & 507{,}450 & 2{,}896{,}896 & 348{,}029 & 4{,}182{,}400 & 267{,}407 & 5{,}185{,}280 & 213{,}703 & 5{,}586{,}304 \\
case\_019 & 81{,}144 & 43{,}224 & 27{,}168 & 29{,}062 & 27{,}168 & 24{,}031 & 256{,}128 & 19{,}200 & 94{,}848 \\
case\_020 & 532{,}746 & 266{,}580 & 0 & 177{,}950 & 0 & 133{,}704 & 0 & 107{,}650 & 0 \\
case\_021 & 6{,}567{,}937 & 3{,}112{,}600 & 8{,}908{,}800 & 2{,}096{,}552 & 12{,}902{,}400 & 1{,}550{,}324 & 4{,}608{,}000 & 1{,}243{,}902 & 13{,}516{,}800 \\
case\_022 & 56{,}780 & 28{,}333 & 6{,}912 & 19{,}053 & 13{,}824 & 14{,}654 & 13{,}824 & 11{,}828 & 20{,}736 \\
case\_023 & 122{,}614 & 63{,}076 & 122{,}936 & 41{,}662 & 196{,}896 & 31{,}083 & 281{,}280 & 24{,}947 & 270{,}912 \\
case\_024 & 2{,}555{,}343 & 2{,}544{,}732 & 393{,}264 & 2{,}267{,}368 & 39{,}329{,}468 & 1{,}799{,}130 & 39{,}328{,}684 & 1{,}799{,}130 & 39{,}328{,}684 \\
case\_025 & 4{,}863{,}026 & 2{,}431{,}840 & 27{,}648 & 1{,}626{,}688 & 552{,}960 & 1{,}596{,}420 & 57{,}474{,}048 & 1{,}364{,}832 & 57{,}830{,}400 \\
case\_026 & 99{,}203 & 62{,}589 & 606{,}260 & 41{,}187 & 217{,}872 & 41{,}187 & 217{,}872 & 41{,}187 & 217{,}872 \\
case\_027 & 2{,}363{,}952 & 1{,}116{,}984 & 1{,}843{,}200 & 837{,}049 & 3{,}686{,}400 & 559{,}196 & 3{,}317{,}760 & 558{,}753 & 5{,}529{,}600 \\
case\_028 & 43{,}836{,}040 & 20{,}572{,}367 & 48{,}637{,}056 & 13{,}753{,}608 & 48{,}529{,}152 & 10{,}408{,}288 & 48{,}817{,}536 & 8{,}372{,}068 & 50{,}895{,}360 \\
case\_029 & 37{,}555 & 19{,}740 & 0 & 12{,}743 & 0 & 11{,}028 & 0 & 9{,}383 & 0 \\
case\_030 & 2{,}808{,}034 & 1{,}396{,}845 & 9{,}431{,}040 & 932{,}172 & 9{,}679{,}360 & 698{,}521 & 7{,}334{,}912 & 560{,}819 & 10{,}660{,}864 \\
case\_031 & 1{,}414{,}369 & 733{,}667 & 6{,}908{,}928 & 485{,}282 & 4{,}462{,}080 & 390{,}234 & 8{,}598{,}820 & 307{,}352 & 8{,}056{,}320 \\
case\_032 & 41{,}443 & 21{,}771 & 18{,}504 & 15{,}388 & 18{,}504 & 14{,}942 & 92{,}316 & 10{,}220 & 24{,}648 \\
case\_033 & 1{,}774{,}196 & 886{,}462 & 6{,}188{,}064 & 601{,}413 & 9{,}595{,}488 & 453{,}651 & 9{,}513{,}600 & 367{,}551 & 8{,}062{,}080 \\
case\_034 & 842{,}996 & 417{,}250 & 1{,}620{,}992 & 276{,}682 & 2{,}939{,}392 & 212{,}296 & 3{,}256{,}832 & 177{,}825 & 3{,}567{,}616 \\
case\_035 & 169{,}571 & 95{,}406 & 9{,}216 & 95{,}406 & 9{,}216 & 95{,}406 & 9{,}216 & 95{,}406 & 9{,}216 \\
case\_036 & 931{,}510 & 466{,}068 & 1{,}152 & 458{,}829 & 8{,}751{,}360 & 233{,}584 & 3{,}456 & 187{,}182 & 4{,}608 \\
case\_037 & 210{,}001 & 105{,}374 & 204{,}800 & 74{,}038 & 1{,}064{,}960 & 53{,}320 & 614{,}400 & 43{,}800 & 1{,}925{,}120 \\
case\_038 & 1{,}364{,}555 & 679{,}403 & 4{,}979{,}712 & 455{,}753 & 6{,}337{,}536 & 344{,}959 & 7{,}661{,}056 & 286{,}350 & 8{,}475{,}648 \\
case\_039 & 1{,}435{,}262 & 692{,}660 & 13{,}541{,}376 & 480{,}930 & 15{,}622{,}144 & 395{,}323 & 17{,}604{,}608 & 395{,}323 & 17{,}604{,}608 \\
case\_040 & 253{,}454 & 126{,}251 & 115{,}200 & 110{,}105 & 316{,}416 & 110{,}105 & 316{,}416 & 110{,}105 & 545{,}800 \\
case\_041 & 6{,}003{,}957 & 2{,}970{,}011 & 36{,}372{,}480 & 1{,}984{,}030 & 40{,}424{,}448 & 1{,}490{,}758 & 36{,}624{,}384 & 1{,}196{,}239 & 39{,}533{,}568 \\
case\_042 & 1{,}686{,}740 & 842{,}181 & 6{,}817{,}056 & 584{,}180 & 8{,}561{,}952 & 433{,}784 & 10{,}958{,}976 & 359{,}815 & 11{,}529{,}504 \\
case\_043 & 677{,}656 & 340{,}574 & 737{,}280 & 268{,}448 & 2{,}257{,}920 & 268{,}172 & 2{,}398{,}342 & 268{,}172 & 2{,}398{,}342 \\
case\_044 & 154{,}407 & 61{,}545 & 1{,}395{,}040 & 61{,}545 & 1{,}395{,}040 & 61{,}545 & 1{,}395{,}040 & 61{,}545 & 1{,}395{,}040 \\
case\_045 & 125{,}614 & 69{,}275 & 368{,}640 & 45{,}572 & 0 & 33{,}246 & 0 & 29{,}135 & 0 \\
case\_046 & 276{,}455 & 128{,}056 & 1{,}192{,}544 & 102{,}504 & 2{,}321{,}600 & 102{,}504 & 2{,}321{,}600 & 102{,}504 & 2{,}321{,}600 \\
case\_047 & 326{,}508 & 237{,}888 & 2{,}548{,}464 & 203{,}634 & 2{,}716{,}176 & 186{,}124 & 2{,}918{,}704 & 178{,}918 & 3{,}102{,}288 \\
case\_048 & 222{,}220 & 149{,}315 & 1{,}043{,}010 & 135{,}083 & 1{,}107{,}632 & 135{,}083 & 1{,}107{,}632 & 135{,}083 & 1{,}107{,}632 \\
case\_049 & 196{,}276 & 147{,}531 & 2{,}868{,}130 & 131{,}522 & 2{,}981{,}982 & 122{,}133 & 2{,}962{,}918 & 122{,}133 & 2{,}962{,}918 \\
case\_050 & 149{,}690 & 106{,}870 & 74{,}860 & 96{,}346 & 2{,}723{,}246 & 89{,}427 & 2{,}725{,}404 & 89{,}427 & 2{,}725{,}404 \\
case\_051 & 607{,}628 & 558{,}663 & 21{,}627{,}306 & 436{,}197 & 11{,}273{,}748 & 422{,}863 & 9{,}045{,}668 & 422{,}863 & 9{,}045{,}668 \\
case\_052 & 180{,}301 & 101{,}771 & 2{,}663{,}284 & 81{,}151 & 2{,}057{,}578 & 61{,}662 & 407{,}040 & 50{,}807 & 443{,}904 \\
case\_053 & 1{,}472{,}686 & 725{,}195 & 10{,}517{,}152 & 539{,}122 & 6{,}888{,}192 & 492{,}232 & 11{,}271{,}680 & 484{,}927 & 11{,}937{,}528 \\
case\_054 & 815{,}051 & 517{,}960 & 1{,}130{,}496 & 457{,}565 & 11{,}226{,}296 & 405{,}511 & 11{,}493{,}906 & 376{,}524 & 11{,}981{,}934 \\
case\_055 & 118{,}835 & 56{,}840 & 99{,}328 & 38{,}542 & 115{,}712 & 35{,}590 & 243{,}724 & 26{,}469 & 198{,}656 \\
case\_056 & 277{,}285 & 215{,}107 & 2{,}685{,}186 & 186{,}802 & 3{,}021{,}450 & 175{,}720 & 3{,}035{,}580 & 166{,}663 & 3{,}074{,}222 \\
case\_057 & 59{,}726 & 33{,}702 & 93{,}988 & 25{,}250 & 14{,}112 & 20{,}957 & 14{,}112 & 16{,}502 & 69{,}696 \\
case\_058 & 4{,}539{,}717 & 2{,}193{,}367 & 38{,}522{,}880 & 1{,}526{,}970 & 56{,}279{,}040 & 1{,}145{,}057 & 47{,}173{,}632 & 1{,}004{,}819 & 32{,}104{,}448 \\
case\_059 & 829{,}473 & 369{,}318 & 360{,}448 & 255{,}117 & 6{,}352{,}896 & 185{,}216 & 1{,}081{,}344 & 150{,}736 & 4{,}866{,}048 \\
case\_060 & 888{,}789 & 471{,}857 & 3{,}890{,}320 & 340{,}707 & 3{,}736{,}840 & 291{,}315 & 3{,}973{,}392 & 249{,}327 & 4{,}516{,}232 \\
case\_061 & 77{,}191 & 55{,}396 & 1{,}162{,}876 & 43{,}481 & 1{,}265{,}444 & 23{,}455 & 64{,}512 & 22{,}283 & 182{,}784 \\
case\_062 & 2{,}854{,}809 & 1{,}417{,}218 & 58{,}354{,}176 & 1{,}079{,}181 & 10{,}572{,}288 & 1{,}079{,}181 & 10{,}572{,}288 & 1{,}079{,}181 & 10{,}572{,}288 \\
case\_063 & 1{,}070{,}397 & 529{,}854 & 21{,}474{,}816 & 413{,}784 & 4{,}399{,}104 & 316{,}530 & 3{,}806{,}208 & 259{,}997 & 3{,}757{,}056 \\
case\_064 & 19{,}116 & 16{,}926 & 81{,}696 & 16{,}926 & 81{,}696 & 16{,}926 & 81{,}696 & 16{,}926 & 81{,}696 \\
case\_065 & 53{,}068 & 26{,}233 & 14{,}112 & 18{,}001 & 27{,}936 & 14{,}178 & 41{,}904 & 12{,}169 & 55{,}584 \\
case\_066 & 359{,}070 & 207{,}194 & 3{,}560{,}460 & 154{,}448 & 3{,}843{,}658 & 143{,}175 & 4{,}313{,}164 & 132{,}051 & 4{,}607{,}842 \\
case\_067 & 63{,}834{,}834 & 30{,}462{,}883 & 74{,}951{,}040 & 20{,}430{,}262 & 77{,}820{,}672 & 15{,}482{,}490 & 81{,}542{,}272 & 12{,}751{,}707 & 90{,}050{,}048 \\
case\_068 & 362{,}095 & 227{,}799 & 2{,}740{,}574 & 187{,}938 & 3{,}674{,}512 & 169{,}482 & 3{,}735{,}730 & 168{,}776 & 3{,}808{,}662 \\
case\_069 & 25{,}144 & 25{,}111 & 68{,}740 & 21{,}307 & 426{,}484 & 19{,}837 & 386{,}908 & 19{,}038 & 552{,}942 \\
case\_070 & 117{,}300 & 70{,}357 & 700{,}584 & 41{,}299 & 273{,}408 & 41{,}299 & 273{,}408 & 30{,}598 & 462{,}848 \\
case\_071 & 18{,}919 & 18{,}919 & 0 & 15{,}923 & 350{,}698 & 14{,}867 & 353{,}802 & 14{,}757 & 343{,}264 \\
case\_072 & 23{,}897{,}828 & 11{,}619{,}091 & 116{,}910{,}880 & 7{,}842{,}498 & 121{,}621{,}280 & 5{,}885{,}034 & 124{,}709{,}664 & 4{,}788{,}636 & 129{,}778{,}352 \\
case\_073 & 10{,}081{,}626 & 4{,}344{,}547 & 16{,}800{,}640 & 2{,}967{,}238 & 19{,}443{,}968 & 2{,}356{,}726 & 20{,}829{,}312 & 2{,}005{,}012 & 24{,}550{,}912 \\
case\_074 & 1{,}682{,}239 & 817{,}481 & 7{,}168{,}000 & 545{,}581 & 8{,}437{,}760 & 437{,}729 & 13{,}484{,}036 & 335{,}235 & 11{,}317{,}248 \\
case\_075 & 1{,}441{,}809 & 824{,}894 & 8{,}622{,}006 & 686{,}606 & 10{,}055{,}794 & 507{,}973 & 10{,}511{,}278 & 465{,}714 & 11{,}274{,}346 \\
case\_076 & 6{,}185{,}113 & 3{,}044{,}060 & 42{,}184{,}704 & 2{,}034{,}436 & 42{,}454{,}016 & 1{,}523{,}743 & 42{,}636{,}288 & 1{,}228{,}370 & 36{,}692{,}992 \\
case\_077 & 495{,}320 & 249{,}892 & 57{,}344 & 177{,}805 & 114{,}688 & 145{,}044 & 172{,}032 & 120{,}238 & 172{,}032 \\
case\_078 & 2{,}296{,}143 & 1{,}080{,}972 & 15{,}039{,}360 & 840{,}164 & 6{,}508{,}800 & 619{,}540 & 15{,}039{,}360 & 605{,}119 & 10{,}578{,}432 \\
case\_079 & 31{,}052{,}251 & 14{,}734{,}808 & 40{,}415{,}232 & 9{,}817{,}862 & 44{,}064{,}768 & 7{,}392{,}284 & 47{,}714{,}304 & 5{,}933{,}675 & 59{,}271{,}168 \\
case\_080 & 292{,}138 & 145{,}926 & 2{,}191{,}360 & 100{,}325 & 3{,}047{,}424 & 100{,}325 & 3{,}047{,}424 & 100{,}325 & 3{,}047{,}424 \\
case\_081 & 151{,}795 & 76{,}860 & 0 & 51{,}943 & 0 & 39{,}588 & 0 & 30{,}943 & 0 \\
case\_082 & 721{,}041 & 442{,}674 & 4{,}125{,}062 & 337{,}233 & 4{,}181{,}712 & 297{,}594 & 4{,}601{,}250 & 286{,}462 & 5{,}294{,}680 \\
case\_083 & 1{,}057{,}993 & 465{,}508 & 4{,}665{,}952 & 326{,}087 & 4{,}613{,}312 & 271{,}555 & 4{,}232{,}992 & 244{,}647 & 5{,}718{,}400 \\
case\_084 & 2{,}507{,}412 & 1{,}142{,}035 & 40{,}060{,}032 & 836{,}220 & 0 & 696{,}677 & 25{,}121{,}472 & 558{,}010 & 19{,}249{,}920 \\
case\_085 & 3{,}187{,}899 & 2{,}123{,}585 & 27{,}680{,}522 & 1{,}583{,}132 & 40{,}486{,}880 & 1{,}283{,}353 & 40{,}472{,}214 & 1{,}207{,}853 & 32{,}596{,}780 \\
case\_086 & 118{,}403 & 90{,}160 & 3{,}082{,}684 & 83{,}941 & 3{,}316{,}920 & 83{,}368 & 3{,}540{,}042 & 83{,}368 & 3{,}540{,}042 \\
case\_087 & 3{,}974{,}071 & 1{,}964{,}553 & 10{,}667{,}520 & 1{,}373{,}717 & 20{,}459{,}520 & 1{,}069{,}520 & 16{,}911{,}360 & 886{,}289 & 15{,}275{,}520 \\
case\_088 & 243{,}623 & 153{,}452 & 2{,}350{,}620 & 128{,}771 & 3{,}116{,}874 & 115{,}345 & 3{,}053{,}952 & 114{,}724 & 3{,}186{,}920 \\
case\_089 & 693{,}092 & 346{,}832 & 238{,}720 & 251{,}199 & 4{,}268{,}288 & 178{,}160 & 716{,}160 & 144{,}242 & 1{,}331{,}712 \\
case\_090 & 1{,}178{,}363 & 511{,}971 & 2{,}987{,}136 & 391{,}377 & 8{,}410{,}224 & 333{,}771 & 7{,}284{,}096 & 324{,}790 & 10{,}947{,}816 \\
case\_091 & 9{,}186{,}769 & 4{,}537{,}069 & 51{,}184{,}096 & 3{,}037{,}726 & 61{,}616{,}960 & 2{,}294{,}489 & 63{,}015{,}072 & 1{,}864{,}035 & 68{,}942{,}848 \\
case\_092 & 4{,}815{,}554 & 2{,}078{,}533 & 19{,}815{,}008 & 1{,}409{,}413 & 18{,}883{,}776 & 1{,}084{,}563 & 15{,}357{,}728 & 894{,}785 & 18{,}565{,}504 \\
case\_093 & 74{,}205 & 50{,}469 & 1{,}130{,}496 & 25{,}103 & 16{,}384 & 19{,}779 & 24{,}576 & 15{,}876 & 32{,}768 \\
case\_094 & 94{,}892 & 56{,}807 & 1{,}117{,}732 & 46{,}304 & 1{,}535{,}048 & 27{,}338 & 318{,}816 & 25{,}370 & 259{,}200 \\
case\_095 & 2{,}084{,}853 & 1{,}042{,}890 & 0 & 695{,}775 & 0 & 524{,}422 & 0 & 524{,}422 & 0 \\
case\_096 & 84{,}194 & 50{,}311 & 288{,}780 & 34{,}318 & 331{,}084 & 26{,}917 & 529{,}920 & 26{,}917 & 529{,}920 \\
case\_097 & 11{,}491{,}509 & 5{,}161{,}740 & 14{,}483{,}456 & 3{,}545{,}477 & 21{,}192{,}704 & 2{,}591{,}900 & 19{,}595{,}264 & 2{,}097{,}000 & 23{,}429{,}120 \\
case\_098 & 177{,}430 & 89{,}551 & 0 & 89{,}551 & 0 & 89{,}518 & 2{,}112 & 89{,}518 & 2{,}112 \\
case\_099 & 1{,}622{,}999 & 803{,}196 & 6{,}967{,}296 & 536{,}653 & 6{,}381{,}568 & 404{,}781 & 9{,}990{,}144 & 336{,}826 & 10{,}371{,}072 \\
case\_100 & 151{,}223 & 94{,}416 & 667{,}064 & 57{,}549 & 538{,}624 & 57{,}549 & 542{,}720 & 57{,}474 & 542{,}720 \\
\end{longtable}
\end{footnotesize}


\section{问题二逐用例结果}
\begin{footnotesize}
\setlength{\tabcolsep}{2.6pt}
\begin{longtable}{l r rr rr rr rr}
\caption{问题二逐用例结果：Makespan（周期）与总额外搬运量（字节）}\label{tab:app-p2}\\
\toprule
\multirow{2}{*}{用例} & \multirow{2}{*}{单核基准} & \multicolumn{2}{c}{2 核} & \multicolumn{2}{c}{3 核} & \multicolumn{2}{c}{4 核} & \multicolumn{2}{c}{5 核} \\
\cmidrule(lr){3-4}\cmidrule(lr){5-6}\cmidrule(lr){7-8}\cmidrule(lr){9-10}
 & & Makespan & 额外搬运 & Makespan & 额外搬运 & Makespan & 额外搬运 & Makespan & 额外搬运 \\
\midrule
\endfirsthead
\toprule
用例 & 单核基准 & \multicolumn{2}{c}{2 核} & \multicolumn{2}{c}{3 核} & \multicolumn{2}{c}{4 核} & \multicolumn{2}{c}{5 核} \\
\midrule
\endhead
\bottomrule
\endfoot
case\_001 & 233{,}110 & 116{,}868 & 1{,}152 & 78{,}545 & 2{,}304 & 58{,}984 & 3{,}456 & 47{,}502 & 4{,}608 \\
case\_002 & 261{,}945 & 125{,}189 & 3{,}592{,}704 & 86{,}670 & 688{,}128 & 71{,}628 & 749{,}568 & 57{,}751 & 734{,}208 \\
case\_003 & 602{,}212 & 319{,}970 & 4{,}021{,}614 & 218{,}628 & 5{,}586{,}704 & 208{,}603 & 5{,}475{,}838 & 187{,}305 & 6{,}918{,}716 \\
case\_004 & 392{,}357 & 199{,}074 & 110{,}592 & 134{,}043 & 221{,}184 & 104{,}188 & 331{,}776 & 86{,}588 & 2{,}147{,}328 \\
case\_005 & 95{,}618 & 52{,}406 & 1{,}066{,}694 & 45{,}808 & 1{,}778{,}986 & 45{,}808 & 1{,}778{,}986 & 45{,}808 & 1{,}778{,}986 \\
case\_006 & 82{,}837 & 37{,}269 & 31{,}104 & 31{,}151 & 30{,}080 & 31{,}151 & 30{,}080 & 31{,}151 & 30{,}080 \\
case\_007 & 1{,}012{,}562 & 507{,}494 & 27{,}648 & 338{,}094 & 55{,}296 & 254{,}398 & 82{,}944 & 203{,}834 & 110{,}592 \\
case\_008 & 487{,}605 & 244{,}266 & 0 & 163{,}359 & 0 & 123{,}060 & 0 & 100{,}603 & 0 \\
case\_009 & 153{,}261 & 92{,}198 & 0 & 91{,}206 & 1{,}747{,}924 & 84{,}566 & 1{,}874{,}932 & 84{,}566 & 1{,}874{,}932 \\
case\_010 & 86{,}969 & 45{,}272 & 48{,}384 & 30{,}199 & 48{,}384 & 29{,}490 & 716{,}200 & 25{,}107 & 723{,}692 \\
case\_011 & 184{,}697 & 92{,}646 & 180{,}224 & 69{,}625 & 1{,}081{,}344 & 46{,}880 & 540{,}672 & 46{,}880 & 540{,}672 \\
case\_012 & 59{,}088 & 29{,}048 & 28{,}800 & 21{,}594 & 20{,}160 & 16{,}456 & 39{,}168 & 13{,}912 & 49{,}536 \\
case\_013 & 263{,}949 & 132{,}198 & 8{,}192 & 88{,}351 & 16{,}384 & 66{,}982 & 24{,}576 & 54{,}363 & 32{,}768 \\
case\_014 & 17{,}698{,}626 & 8{,}894{,}908 & 66{,}855{,}864 & 5{,}987{,}197 & 56{,}620{,}200 & 4{,}538{,}463 & 80{,}207{,}256 & 3{,}684{,}398 & 64{,}766{,}472 \\
case\_015 & 177{,}362 & 91{,}863 & 109{,}824 & 60{,}183 & 208{,}128 & 49{,}608 & 418{,}572 & 41{,}485 & 317{,}952 \\
case\_016 & 7{,}715{,}523 & 7{,}066{,}708 & 231{,}832 & 3{,}206{,}442 & 18{,}280 & 2{,}548{,}546 & 23{,}168 & 2{,}547{,}814 & 352{,}064 \\
case\_017 & 162{,}810 & 70{,}999 & 87{,}552 & 54{,}520 & 175{,}104 & 39{,}231 & 259{,}200 & 33{,}790 & 340{,}992 \\
case\_018 & 1{,}004{,}367 & 493{,}585 & 1{,}893{,}376 & 341{,}934 & 2{,}675{,}072 & 266{,}513 & 4{,}611{,}840 & 213{,}703 & 5{,}038{,}080 \\
case\_019 & 81{,}144 & 39{,}773 & 37{,}920 & 28{,}658 & 24{,}576 & 23{,}450 & 680{,}776 & 18{,}226 & 67{,}680 \\
case\_020 & 532{,}746 & 266{,}580 & 0 & 177{,}950 & 0 & 133{,}704 & 0 & 107{,}650 & 0 \\
case\_021 & 6{,}567{,}937 & 3{,}194{,}570 & 3{,}174{,}400 & 2{,}144{,}024 & 7{,}766{,}016 & 1{,}550{,}324 & 4{,}608{,}000 & 1{,}250{,}168 & 11{,}431{,}936 \\
case\_022 & 56{,}780 & 28{,}247 & 6{,}912 & 18{,}844 & 13{,}824 & 14{,}555 & 13{,}824 & 11{,}828 & 20{,}736 \\
case\_023 & 122{,}614 & 58{,}256 & 1{,}192{,}892 & 40{,}897 & 196{,}896 & 30{,}755 & 281{,}280 & 24{,}947 & 270{,}912 \\
case\_024 & 2{,}555{,}343 & 2{,}334{,}112 & 230{,}200 & 1{,}062{,}198 & 6{,}040 & 846{,}529 & 7{,}652 & 846{,}529 & 7{,}652 \\
case\_025 & 4{,}863{,}026 & 2{,}431{,}840 & 27{,}648 & 1{,}622{,}164 & 55{,}296 & 1{,}217{,}444 & 82{,}944 & 975{,}282 & 110{,}592 \\
case\_026 & 99{,}203 & 47{,}565 & 1{,}296 & 41{,}536 & 95{,}664 & 41{,}536 & 95{,}664 & 38{,}400 & 1{,}162{,}436 \\
case\_027 & 2{,}363{,}952 & 1{,}120{,}962 & 1{,}744{,}896 & 837{,}049 & 3{,}686{,}400 & 559{,}196 & 3{,}317{,}760 & 558{,}753 & 5{,}529{,}600 \\
case\_028 & 43{,}836{,}040 & 20{,}475{,}783 & 25{,}833{,}600 & 13{,}590{,}237 & 23{,}420{,}160 & 10{,}268{,}731 & 36{,}349{,}824 & 8{,}169{,}383 & 33{,}460{,}224 \\
case\_029 & 37{,}555 & 19{,}740 & 0 & 12{,}743 & 0 & 11{,}028 & 0 & 9{,}383 & 0 \\
case\_030 & 2{,}808{,}034 & 1{,}390{,}756 & 2{,}045{,}952 & 937{,}837 & 8{,}459{,}264 & 696{,}635 & 4{,}763{,}648 & 560{,}819 & 10{,}660{,}864 \\
case\_031 & 1{,}414{,}369 & 734{,}985 & 3{,}924{,}480 & 486{,}170 & 4{,}591{,}104 & 390{,}627 & 5{,}976{,}868 & 307{,}985 & 7{,}549{,}440 \\
case\_032 & 41{,}443 & 19{,}737 & 6{,}216 & 12{,}854 & 12{,}360 & 12{,}854 & 12{,}360 & 10{,}415 & 12{,}360 \\
case\_033 & 1{,}774{,}196 & 887{,}381 & 6{,}501{,}408 & 602{,}654 & 6{,}600{,}768 & 459{,}481 & 7{,}844{,}448 & 376{,}887 & 7{,}767{,}680 \\
case\_034 & 842{,}996 & 412{,}835 & 1{,}082{,}880 & 271{,}887 & 2{,}453{,}504 & 207{,}388 & 3{,}150{,}336 & 170{,}400 & 3{,}487{,}744 \\
case\_035 & 169{,}571 & 95{,}406 & 0 & 95{,}406 & 0 & 95{,}406 & 0 & 95{,}406 & 0 \\
case\_036 & 931{,}510 & 466{,}068 & 1{,}152 & 311{,}345 & 2{,}304 & 233{,}584 & 3{,}456 & 187{,}182 & 4{,}608 \\
case\_037 & 210{,}001 & 105{,}374 & 204{,}800 & 74{,}038 & 1{,}064{,}960 & 53{,}320 & 614{,}400 & 43{,}800 & 1{,}925{,}120 \\
case\_038 & 1{,}364{,}555 & 678{,}569 & 1{,}992{,}192 & 456{,}517 & 5{,}171{,}712 & 346{,}741 & 8{,}735{,}232 & 285{,}550 & 8{,}265{,}216 \\
case\_039 & 1{,}435{,}262 & 708{,}561 & 11{,}067{,}392 & 485{,}565 & 14{,}323{,}712 & 395{,}323 & 17{,}604{,}608 & 395{,}323 & 17{,}604{,}608 \\
case\_040 & 253{,}454 & 120{,}986 & 12{,}288 & 100{,}177 & 65{,}090 & 99{,}511 & 265{,}034 & 95{,}285 & 388{,}940 \\
case\_041 & 6{,}003{,}957 & 2{,}989{,}434 & 26{,}674{,}176 & 2{,}006{,}212 & 33{,}024{,}000 & 1{,}497{,}746 & 34{,}397{,}184 & 1{,}203{,}317 & 36{,}360{,}192 \\
case\_042 & 1{,}686{,}740 & 845{,}035 & 6{,}009{,}120 & 584{,}180 & 8{,}561{,}952 & 438{,}568 & 8{,}294{,}688 & 355{,}113 & 10{,}383{,}360 \\
case\_043 & 677{,}656 & 338{,}088 & 460{,}800 & 245{,}028 & 352{,}450 & 245{,}028 & 352{,}450 & 233{,}701 & 7{,}587{,}124 \\
case\_044 & 154{,}407 & 43{,}931 & 930{,}400 & 43{,}931 & 930{,}400 & 43{,}931 & 930{,}400 & 43{,}931 & 930{,}400 \\
case\_045 & 125{,}614 & 63{,}059 & 0 & 42{,}601 & 0 & 31{,}878 & 0 & 26{,}164 & 0 \\
case\_046 & 276{,}455 & 106{,}042 & 930{,}400 & 80{,}028 & 1{,}860{,}800 & 78{,}210 & 2{,}791{,}200 & 78{,}210 & 2{,}791{,}200 \\
case\_047 & 326{,}508 & 177{,}730 & 2{,}253{,}222 & 138{,}079 & 3{,}288{,}588 & 128{,}508 & 4{,}357{,}842 & 128{,}508 & 4{,}357{,}842 \\
case\_048 & 222{,}220 & 129{,}426 & 314{,}182 & 113{,}494 & 456{,}716 & 111{,}484 & 1{,}039{,}728 & 111{,}484 & 1{,}039{,}728 \\
case\_049 & 196{,}276 & 126{,}134 & 1{,}316{,}902 & 99{,}627 & 1{,}901{,}512 & 89{,}606 & 2{,}249{,}534 & 89{,}606 & 2{,}249{,}534 \\
case\_050 & 149{,}690 & 84{,}490 & 854{,}586 & 72{,}131 & 1{,}528{,}238 & 64{,}647 & 1{,}809{,}716 & 64{,}647 & 1{,}809{,}716 \\
case\_051 & 607{,}628 & 547{,}789 & 229{,}584 & 252{,}851 & 1{,}420 & 202{,}281 & 1{,}724 & 201{,}161 & 1{,}808 \\
case\_052 & 180{,}301 & 92{,}135 & 1{,}938{,}984 & 68{,}329 & 1{,}409{,}612 & 61{,}312 & 259{,}584 & 50{,}551 & 328{,}704 \\
case\_053 & 1{,}472{,}686 & 698{,}368 & 8{,}085{,}394 & 506{,}815 & 5{,}882{,}944 & 477{,}534 & 6{,}416{,}288 & 443{,}726 & 15{,}484{,}744 \\
case\_054 & 815{,}051 & 472{,}156 & 6{,}564{,}694 & 347{,}229 & 9{,}468{,}104 & 317{,}609 & 7{,}167{,}478 & 287{,}243 & 11{,}160{,}124 \\
case\_055 & 118{,}835 & 56{,}668 & 57{,}856 & 38{,}406 & 115{,}712 & 31{,}367 & 115{,}712 & 26{,}322 & 198{,}656 \\
case\_056 & 277{,}285 & 172{,}161 & 1{,}922{,}368 & 129{,}190 & 2{,}780{,}198 & 124{,}871 & 3{,}387{,}698 & 124{,}871 & 3{,}387{,}698 \\
case\_057 & 59{,}726 & 33{,}603 & 288 & 25{,}250 & 288 & 20{,}957 & 288 & 16{,}538 & 27{,}936 \\
case\_058 & 4{,}539{,}717 & 2{,}271{,}316 & 33{,}431{,}552 & 1{,}554{,}831 & 30{,}400{,}512 & 1{,}154{,}519 & 42{,}622{,}976 & 1{,}004{,}819 & 32{,}104{,}448 \\
case\_059 & 829{,}473 & 369{,}318 & 360{,}448 & 266{,}217 & 2{,}998{,}272 & 185{,}216 & 1{,}081{,}344 & 150{,}782 & 4{,}325{,}376 \\
case\_060 & 888{,}789 & 464{,}279 & 1{,}705{,}984 & 334{,}694 & 2{,}796{,}672 & 282{,}484 & 3{,}623{,}168 & 237{,}234 & 4{,}286{,}848 \\
case\_061 & 77{,}191 & 47{,}568 & 671{,}274 & 28{,}409 & 43{,}008 & 22{,}111 & 62{,}208 & 18{,}783 & 103{,}936 \\
case\_062 & 2{,}854{,}809 & 1{,}718{,}716 & 62{,}917{,}632 & 1{,}056{,}247 & 8{,}732{,}160 & 1{,}056{,}247 & 8{,}732{,}160 & 1{,}056{,}247 & 8{,}732{,}160 \\
case\_063 & 1{,}070{,}397 & 618{,}185 & 20{,}963{,}328 & 397{,}260 & 3{,}081{,}216 & 301{,}177 & 2{,}672{,}640 & 244{,}345 & 2{,}939{,}904 \\
case\_064 & 19{,}116 & 13{,}414 & 77{,}222 & 12{,}280 & 94{,}666 & 11{,}941 & 143{,}022 & 11{,}406 & 142{,}706 \\
case\_065 & 53{,}068 & 25{,}997 & 13{,}968 & 18{,}001 & 27{,}936 & 14{,}178 & 41{,}904 & 12{,}135 & 55{,}584 \\
case\_066 & 359{,}070 & 183{,}389 & 1{,}932{,}976 & 128{,}156 & 2{,}549{,}320 & 110{,}097 & 3{,}362{,}510 & 109{,}083 & 3{,}836{,}102 \\
case\_067 & 63{,}834{,}834 & 30{,}476{,}248 & 56{,}984{,}192 & 20{,}333{,}946 & 17{,}584{,}896 & 15{,}271{,}913 & 16{,}931{,}968 & 12{,}621{,}657 & 17{,}327{,}616 \\
case\_068 & 362{,}095 & 202{,}721 & 2{,}254{,}270 & 157{,}301 & 3{,}594{,}672 & 143{,}171 & 4{,}306{,}378 & 143{,}025 & 3{,}353{,}518 \\
case\_069 & 25{,}144 & 13{,}965 & 404{,}806 & 13{,}965 & 404{,}806 & 13{,}965 & 404{,}806 & 13{,}831 & 377{,}328 \\
case\_070 & 117{,}300 & 56{,}421 & 129{,}536 & 40{,}990 & 259{,}072 & 33{,}697 & 360{,}960 & 30{,}598 & 462{,}848 \\
case\_071 & 18{,}919 & 11{,}580 & 269{,}862 & 11{,}188 & 400{,}138 & 11{,}188 & 400{,}138 & 11{,}188 & 400{,}138 \\
case\_072 & 23{,}897{,}828 & 11{,}899{,}903 & 91{,}122{,}464 & 8{,}008{,}124 & 113{,}388{,}320 & 5{,}993{,}260 & 122{,}669{,}856 & 4{,}861{,}592 & 113{,}066{,}672 \\
case\_073 & 10{,}081{,}626 & 4{,}238{,}435 & 3{,}721{,}600 & 2{,}825{,}285 & 7{,}443{,}200 & 2{,}225{,}291 & 11{,}164{,}800 & 1{,}849{,}407 & 14{,}886{,}400 \\
case\_074 & 1{,}682{,}239 & 811{,}260 & 7{,}262{,}208 & 542{,}279 & 7{,}397{,}376 & 409{,}816 & 8{,}658{,}948 & 341{,}700 & 8{,}777{,}728 \\
case\_075 & 1{,}441{,}809 & 689{,}383 & 3{,}963{,}590 & 570{,}824 & 5{,}602{,}700 & 408{,}170 & 6{,}136{,}978 & 377{,}686 & 8{,}326{,}392 \\
case\_076 & 6{,}185{,}113 & 3{,}094{,}537 & 30{,}394{,}368 & 2{,}068{,}186 & 31{,}798{,}272 & 1{,}535{,}669 & 26{,}690{,}560 & 1{,}246{,}135 & 35{,}584{,}000 \\
case\_077 & 495{,}320 & 249{,}722 & 57{,}344 & 177{,}575 & 114{,}688 & 145{,}044 & 172{,}032 & 120{,}068 & 172{,}032 \\
case\_078 & 2{,}296{,}143 & 1{,}024{,}634 & 2{,}148{,}480 & 768{,}769 & 4{,}849{,}920 & 520{,}320 & 6{,}445{,}440 & 520{,}320 & 6{,}445{,}440 \\
case\_079 & 31{,}052{,}251 & 15{,}446{,}486 & 31{,}780{,}864 & 10{,}181{,}692 & 39{,}960{,}576 & 7{,}635{,}624 & 30{,}445{,}568 & 6{,}075{,}976 & 47{,}677{,}440 \\
case\_080 & 292{,}138 & 145{,}926 & 2{,}191{,}360 & 100{,}325 & 3{,}047{,}424 & 100{,}325 & 3{,}047{,}424 & 100{,}325 & 3{,}047{,}424 \\
case\_081 & 151{,}795 & 76{,}860 & 0 & 51{,}943 & 0 & 39{,}588 & 0 & 30{,}943 & 0 \\
case\_082 & 721{,}041 & 366{,}243 & 2{,}233{,}350 & 264{,}390 & 3{,}467{,}596 & 229{,}458 & 4{,}577{,}264 & 204{,}292 & 6{,}027{,}608 \\
case\_083 & 1{,}057{,}993 & 419{,}777 & 930{,}400 & 289{,}139 & 1{,}860{,}800 & 211{,}087 & 2{,}791{,}200 & 185{,}314 & 3{,}721{,}600 \\
case\_084 & 2{,}507{,}412 & 1{,}255{,}052 & 0 & 836{,}220 & 0 & 629{,}028 & 0 & 503{,}472 & 0 \\
case\_085 & 3{,}187{,}899 & 1{,}743{,}831 & 13{,}942{,}986 & 1{,}246{,}519 & 21{,}276{,}524 & 946{,}039 & 21{,}333{,}874 & 846{,}587 & 25{,}456{,}536 \\
case\_086 & 118{,}403 & 65{,}985 & 1{,}576{,}580 & 57{,}154 & 2{,}169{,}644 & 56{,}538 & 2{,}561{,}616 & 56{,}538 & 2{,}561{,}616 \\
case\_087 & 3{,}974{,}071 & 1{,}985{,}234 & 9{,}999{,}360 & 1{,}379{,}480 & 9{,}183{,}744 & 1{,}069{,}661 & 11{,}672{,}064 & 885{,}851 & 8{,}663{,}040 \\
case\_088 & 243{,}623 & 133{,}265 & 1{,}295{,}018 & 105{,}066 & 2{,}175{,}676 & 90{,}485 & 2{,}472{,}754 & 90{,}485 & 2{,}472{,}754 \\
case\_089 & 693{,}092 & 295{,}856 & 238{,}720 & 208{,}395 & 1{,}399{,}040 & 148{,}856 & 716{,}160 & 120{,}170 & 954{,}880 \\
case\_090 & 1{,}178{,}363 & 451{,}762 & 2{,}165{,}760 & 333{,}250 & 4{,}521{,}600 & 306{,}241 & 6{,}670{,}080 & 276{,}562 & 8{,}593{,}920 \\
case\_091 & 9{,}186{,}769 & 4{,}588{,}580 & 42{,}713{,}056 & 3{,}069{,}210 & 46{,}052{,}672 & 2{,}315{,}629 & 66{,}681{,}504 & 1{,}872{,}971 & 53{,}916{,}160 \\
case\_092 & 4{,}815{,}554 & 2{,}012{,}858 & 8{,}012{,}384 & 1{,}352{,}222 & 7{,}431{,}360 & 1{,}016{,}169 & 6{,}264{,}608 & 797{,}596 & 3{,}950{,}976 \\
case\_093 & 74{,}205 & 37{,}827 & 8{,}192 & 25{,}103 & 16{,}384 & 19{,}779 & 24{,}576 & 15{,}876 & 32{,}768 \\
case\_094 & 94{,}892 & 46{,}959 & 106{,}272 & 38{,}877 & 1{,}143{,}872 & 26{,}469 & 318{,}816 & 24{,}751 & 259{,}200 \\
case\_095 & 2{,}084{,}853 & 1{,}042{,}890 & 0 & 695{,}775 & 0 & 524{,}422 & 0 & 420{,}852 & 0 \\
case\_096 & 84{,}194 & 46{,}739 & 140{,}168 & 31{,}518 & 143{,}884 & 26{,}984 & 78{,}336 & 26{,}984 & 78{,}336 \\
case\_097 & 11{,}491{,}509 & 5{,}617{,}602 & 11{,}796{,}480 & 3{,}660{,}459 & 16{,}031{,}744 & 2{,}626{,}908 & 14{,}221{,}312 & 2{,}102{,}348 & 21{,}311{,}488 \\
case\_098 & 177{,}430 & 89{,}551 & 0 & 89{,}551 & 0 & 88{,}241 & 3{,}435{,}400 & 64{,}329 & 2{,}258{,}230 \\
case\_099 & 1{,}622{,}999 & 804{,}729 & 6{,}524{,}928 & 537{,}579 & 6{,}754{,}304 & 405{,}898 & 9{,}334{,}784 & 331{,}823 & 8{,}712{,}192 \\
case\_100 & 151{,}223 & 73{,}847 & 72{,}192 & 57{,}811 & 72{,}192 & 57{,}811 & 72{,}192 & 57{,}811 & 72{,}192 \\
\end{longtable}
\end{footnotesize}


\section{问题三逐用例结果}
表中“加速比”为同核数下只读 Cache 相对无 L2 的加速比（两种配置各自取最优方案），命中率为只读 Cache 配置下的字节命中率。
\begin{footnotesize}
\setlength{\tabcolsep}{3.2pt}
\begin{longtable}{l c rr c rr c}
\caption{问题三逐用例结果：无 L2 与只读 Cache 两种配置的 Makespan（周期）、总额外搬运量（字节）与 Cache 命中率}\label{tab:app-p3}\\
\toprule
用例 & 核数 & 无 L2 Makespan & Cache Makespan & 加速比 & 无 L2 额外搬运 & Cache 额外搬运 & 命中率 \\
\midrule
\endfirsthead
\toprule
用例 & 核数 & 无 L2 Makespan & Cache Makespan & 加速比 & 无 L2 额外搬运 & Cache 额外搬运 & 命中率 \\
\midrule
\endhead
\bottomrule
\endfoot
case\_001 & 1 & 233{,}110 & 233{,}110 & 1.0000 & 0 & 0 & 0.0\% \\
 & 2 & 116{,}868 & 116{,}868 & 1.0000 & 1{,}152 & 1{,}152 & 0.0\% \\
 & 3 & 78{,}545 & 78{,}545 & 1.0000 & 2{,}304 & 2{,}304 & 0.0\% \\
 & 4 & 58{,}984 & 58{,}984 & 1.0000 & 3{,}456 & 3{,}456 & 0.0\% \\
 & 5 & 47{,}502 & 47{,}502 & 1.0000 & 4{,}608 & 4{,}608 & 0.0\% \\
\addlinespace[2pt]
case\_002 & 1 & 255{,}084 & 255{,}084 & 1.0000 & 3{,}551{,}232 & 3{,}551{,}232 & 34.6\% \\
 & 2 & 125{,}189 & 125{,}102 & 1.0007 & 3{,}592{,}704 & 3{,}592{,}704 & 8.9\% \\
 & 3 & 86{,}670 & 86{,}670 & 1.0000 & 688{,}128 & 688{,}128 & 0.0\% \\
 & 4 & 71{,}628 & 71{,}628 & 1.0000 & 749{,}568 & 749{,}568 & 0.0\% \\
 & 5 & 57{,}751 & 57{,}751 & 1.0000 & 734{,}208 & 734{,}208 & 0.0\% \\
\addlinespace[2pt]
case\_003 & 1 & 576{,}079 & 576{,}079 & 1.0000 & 44{,}068 & 44{,}068 & 0.0\% \\
 & 2 & 319{,}970 & 319{,}550 & 1.0013 & 4{,}021{,}614 & 4{,}021{,}614 & 12.0\% \\
 & 3 & 218{,}628 & 215{,}087 & 1.0165 & 5{,}586{,}704 & 5{,}586{,}704 & 19.3\% \\
 & 4 & 208{,}603 & 203{,}107 & 1.0271 & 5{,}475{,}838 & 8{,}088{,}746 & 24.1\% \\
 & 5 & 187{,}305 & 171{,}861 & 1.0899 & 6{,}918{,}716 & 6{,}918{,}716 & 35.8\% \\
\addlinespace[2pt]
case\_004 & 1 & 392{,}357 & 392{,}357 & 1.0000 & 0 & 0 & 0.0\% \\
 & 2 & 199{,}074 & 199{,}074 & 1.0000 & 110{,}592 & 110{,}592 & 0.0\% \\
 & 3 & 134{,}043 & 134{,}043 & 1.0000 & 221{,}184 & 221{,}184 & 0.0\% \\
 & 4 & 104{,}188 & 103{,}934 & 1.0024 & 331{,}776 & 2{,}036{,}736 & 33.2\% \\
 & 5 & 86{,}588 & 85{,}176 & 1.0166 & 2{,}147{,}328 & 2{,}147{,}328 & 29.3\% \\
\addlinespace[2pt]
case\_005 & 1 & 86{,}400 & 86{,}400 & 1.0000 & 0 & 0 & 0.0\% \\
 & 2 & 52{,}406 & 52{,}390 & 1.0003 & 1{,}066{,}694 & 1{,}066{,}694 & 6.1\% \\
 & 3 & 45{,}808 & 44{,}692 & 1.0250 & 1{,}778{,}986 & 1{,}778{,}986 & 21.1\% \\
 & 4 & 45{,}808 & 43{,}684 & 1.0486 & 1{,}778{,}986 & 2{,}094{,}032 & 30.6\% \\
 & 5 & 45{,}808 & 43{,}684 & 1.0486 & 1{,}778{,}986 & 2{,}094{,}032 & 30.6\% \\
\addlinespace[2pt]
case\_006 & 1 & 75{,}883 & 75{,}883 & 1.0000 & 0 & 0 & 0.0\% \\
 & 2 & 37{,}269 & 37{,}269 & 1.0000 & 31{,}104 & 31{,}104 & 5.4\% \\
 & 3 & 31{,}151 & 31{,}151 & 1.0000 & 30{,}080 & 30{,}080 & 5.2\% \\
 & 4 & 31{,}151 & 26{,}495 & 1.1757 & 30{,}080 & 1{,}231{,}528 & 25.5\% \\
 & 5 & 31{,}151 & 26{,}495 & 1.1757 & 30{,}080 & 1{,}231{,}528 & 25.5\% \\
\addlinespace[2pt]
case\_007 & 1 & 1{,}012{,}562 & 1{,}012{,}562 & 1.0000 & 0 & 0 & 0.0\% \\
 & 2 & 507{,}494 & 507{,}494 & 1.0000 & 27{,}648 & 27{,}648 & 0.0\% \\
 & 3 & 338{,}094 & 338{,}094 & 1.0000 & 55{,}296 & 55{,}296 & 0.0\% \\
 & 4 & 254{,}398 & 254{,}398 & 1.0000 & 82{,}944 & 82{,}944 & 0.0\% \\
 & 5 & 203{,}834 & 203{,}834 & 1.0000 & 110{,}592 & 110{,}592 & 0.0\% \\
\addlinespace[2pt]
case\_008 & 1 & 487{,}605 & 487{,}605 & 1.0000 & 0 & 0 & 0.0\% \\
 & 2 & 244{,}266 & 244{,}266 & 1.0000 & 0 & 0 & 0.0\% \\
 & 3 & 163{,}359 & 163{,}359 & 1.0000 & 0 & 0 & 0.0\% \\
 & 4 & 123{,}060 & 123{,}060 & 1.0000 & 0 & 0 & 0.0\% \\
 & 5 & 100{,}603 & 100{,}603 & 1.0000 & 0 & 0 & 0.0\% \\
\addlinespace[2pt]
case\_009 & 1 & 139{,}006 & 139{,}000 & 1.0000 & 427{,}016 & 427{,}016 & 16.9\% \\
 & 2 & 92{,}198 & 92{,}198 & 1.0000 & 0 & 0 & 0.0\% \\
 & 3 & 91{,}206 & 86{,}965 & 1.0488 & 1{,}747{,}924 & 1{,}747{,}924 & 24.8\% \\
 & 4 & 84{,}566 & 79{,}444 & 1.0645 & 1{,}874{,}932 & 1{,}874{,}932 & 26.2\% \\
 & 5 & 84{,}566 & 79{,}338 & 1.0659 & 1{,}874{,}932 & 2{,}269{,}846 & 29.0\% \\
\addlinespace[2pt]
case\_010 & 1 & 86{,}969 & 86{,}969 & 1.0000 & 0 & 0 & 0.0\% \\
 & 2 & 45{,}272 & 45{,}272 & 1.0000 & 48{,}384 & 48{,}384 & 11.2\% \\
 & 3 & 30{,}199 & 30{,}199 & 1.0000 & 48{,}384 & 48{,}384 & 11.2\% \\
 & 4 & 29{,}490 & 28{,}773 & 1.0249 & 716{,}200 & 716{,}200 & 30.6\% \\
 & 5 & 25{,}107 & 24{,}809 & 1.0120 & 723{,}692 & 723{,}692 & 32.1\% \\
\addlinespace[2pt]
case\_011 & 1 & 184{,}697 & 184{,}697 & 1.0000 & 0 & 0 & 0.0\% \\
 & 2 & 92{,}646 & 92{,}646 & 1.0000 & 180{,}224 & 180{,}224 & 0.0\% \\
 & 3 & 69{,}625 & 69{,}625 & 1.0000 & 1{,}081{,}344 & 1{,}081{,}344 & 56.7\% \\
 & 4 & 46{,}880 & 46{,}880 & 1.0000 & 540{,}672 & 540{,}672 & 0.0\% \\
 & 5 & 46{,}880 & 46{,}777 & 1.0022 & 540{,}672 & 1{,}622{,}016 & 62.3\% \\
\addlinespace[2pt]
case\_012 & 1 & 55{,}629 & 55{,}507 & 1.0022 & 67{,}604 & 81{,}940 & 11.1\% \\
 & 2 & 29{,}048 & 29{,}048 & 1.0000 & 28{,}800 & 28{,}800 & 5.6\% \\
 & 3 & 21{,}594 & 21{,}594 & 1.0000 & 20{,}160 & 20{,}160 & 4.0\% \\
 & 4 & 16{,}456 & 16{,}456 & 1.0000 & 39{,}168 & 39{,}168 & 7.5\% \\
 & 5 & 13{,}912 & 13{,}912 & 1.0000 & 49{,}536 & 49{,}536 & 7.0\% \\
\addlinespace[2pt]
case\_013 & 1 & 263{,}949 & 263{,}949 & 1.0000 & 0 & 0 & 0.0\% \\
 & 2 & 132{,}198 & 132{,}198 & 1.0000 & 8{,}192 & 8{,}192 & 0.0\% \\
 & 3 & 88{,}351 & 88{,}351 & 1.0000 & 16{,}384 & 16{,}384 & 0.0\% \\
 & 4 & 66{,}982 & 66{,}982 & 1.0000 & 24{,}576 & 24{,}576 & 0.0\% \\
 & 5 & 54{,}363 & 54{,}363 & 1.0000 & 32{,}768 & 32{,}768 & 0.0\% \\
\addlinespace[2pt]
case\_014 & 1 & 17{,}698{,}626 & 17{,}689{,}602 & 1.0005 & 71{,}010{,}720 & 71{,}010{,}720 & 12.3\% \\
 & 2 & 8{,}894{,}908 & 8{,}882{,}659 & 1.0014 & 66{,}855{,}864 & 66{,}855{,}864 & 31.6\% \\
 & 3 & 5{,}987{,}197 & 5{,}944{,}828 & 1.0071 & 56{,}620{,}200 & 56{,}620{,}200 & 38.7\% \\
 & 4 & 4{,}538{,}463 & 4{,}518{,}920 & 1.0043 & 80{,}207{,}256 & 80{,}207{,}256 & 33.6\% \\
 & 5 & 3{,}684{,}398 & 3{,}670{,}307 & 1.0038 & 64{,}766{,}472 & 64{,}766{,}472 & 44.0\% \\
\addlinespace[2pt]
case\_015 & 1 & 170{,}458 & 170{,}458 & 1.0000 & 24{,}576 & 24{,}576 & 2.7\% \\
 & 2 & 91{,}863 & 91{,}863 & 1.0000 & 109{,}824 & 109{,}824 & 11.1\% \\
 & 3 & 60{,}183 & 60{,}183 & 1.0000 & 208{,}128 & 208{,}128 & 16.1\% \\
 & 4 & 49{,}608 & 49{,}608 & 1.0000 & 418{,}572 & 418{,}572 & 21.1\% \\
 & 5 & 41{,}485 & 41{,}368 & 1.0028 & 317{,}952 & 317{,}952 & 21.0\% \\
\addlinespace[2pt]
case\_016 & 1 & 7{,}715{,}523 & 7{,}715{,}523 & 1.0000 & 0 & 0 & 0.0\% \\
 & 2 & 7{,}066{,}708 & 7{,}066{,}708 & 1.0000 & 231{,}832 & 231{,}832 & 36.8\% \\
 & 3 & 3{,}206{,}442 & 3{,}206{,}442 & 1.0000 & 18{,}280 & 18{,}280 & 0.1\% \\
 & 4 & 2{,}548{,}546 & 2{,}548{,}546 & 1.0000 & 23{,}168 & 23{,}168 & 0.3\% \\
 & 5 & 2{,}547{,}814 & 2{,}547{,}816 & 1.0000 & 352{,}064 & 24{,}384 & 0.4\% \\
\addlinespace[2pt]
case\_017 & 1 & 151{,}167 & 151{,}167 & 1.0000 & 0 & 0 & 0.0\% \\
 & 2 & 70{,}999 & 70{,}997 & 1.0000 & 87{,}552 & 87{,}552 & 14.1\% \\
 & 3 & 54{,}520 & 54{,}462 & 1.0011 & 175{,}104 & 175{,}104 & 3.7\% \\
 & 4 & 39{,}231 & 39{,}231 & 1.0000 & 259{,}200 & 259{,}200 & 17.2\% \\
 & 5 & 33{,}790 & 33{,}790 & 1.0000 & 340{,}992 & 340{,}992 & 4.2\% \\
\addlinespace[2pt]
case\_018 & 1 & 1{,}004{,}367 & 1{,}003{,}110 & 1.0013 & 1{,}970{,}176 & 1{,}970{,}176 & 34.4\% \\
 & 2 & 493{,}585 & 493{,}325 & 1.0005 & 1{,}893{,}376 & 1{,}893{,}376 & 45.9\% \\
 & 3 & 341{,}934 & 341{,}934 & 1.0000 & 2{,}675{,}072 & 2{,}675{,}072 & 60.1\% \\
 & 4 & 266{,}513 & 266{,}431 & 1.0003 & 4{,}611{,}840 & 4{,}611{,}840 & 61.5\% \\
 & 5 & 213{,}703 & 212{,}418 & 1.0060 & 5{,}038{,}080 & 5{,}038{,}080 & 60.0\% \\
\addlinespace[2pt]
case\_019 & 1 & 65{,}302 & 65{,}302 & 1.0000 & 272{,}384 & 272{,}384 & 8.1\% \\
 & 2 & 39{,}773 & 39{,}773 & 1.0000 & 37{,}920 & 37{,}920 & 3.4\% \\
 & 3 & 28{,}658 & 28{,}658 & 1.0000 & 24{,}576 & 24{,}576 & 4.0\% \\
 & 4 & 23{,}450 & 23{,}149 & 1.0130 & 680{,}776 & 680{,}776 & 18.7\% \\
 & 5 & 18{,}226 & 18{,}226 & 1.0000 & 67{,}680 & 67{,}680 & 11.8\% \\
\addlinespace[2pt]
case\_020 & 1 & 532{,}746 & 532{,}746 & 1.0000 & 0 & 0 & 0.0\% \\
 & 2 & 266{,}580 & 266{,}580 & 1.0000 & 0 & 0 & 0.0\% \\
 & 3 & 177{,}950 & 177{,}950 & 1.0000 & 0 & 0 & 0.0\% \\
 & 4 & 133{,}704 & 133{,}704 & 1.0000 & 0 & 0 & 0.0\% \\
 & 5 & 107{,}650 & 107{,}650 & 1.0000 & 0 & 0 & 0.0\% \\
\addlinespace[2pt]
case\_021 & 1 & 6{,}567{,}937 & 6{,}558{,}837 & 1.0014 & 4{,}501{,}504 & 4{,}501{,}504 & 50.1\% \\
 & 2 & 3{,}194{,}570 & 3{,}193{,}634 & 1.0003 & 3{,}174{,}400 & 3{,}174{,}400 & 21.8\% \\
 & 3 & 2{,}144{,}024 & 2{,}142{,}851 & 1.0005 & 7{,}766{,}016 & 7{,}766{,}016 & 54.6\% \\
 & 4 & 1{,}550{,}324 & 1{,}550{,}324 & 1.0000 & 4{,}608{,}000 & 4{,}608{,}000 & 0.0\% \\
 & 5 & 1{,}250{,}168 & 1{,}249{,}380 & 1.0006 & 11{,}431{,}936 & 11{,}431{,}936 & 57.0\% \\
\addlinespace[2pt]
case\_022 & 1 & 56{,}780 & 56{,}780 & 1.0000 & 0 & 0 & 0.0\% \\
 & 2 & 28{,}247 & 28{,}247 & 1.0000 & 6{,}912 & 6{,}912 & 3.1\% \\
 & 3 & 18{,}844 & 18{,}844 & 1.0000 & 13{,}824 & 13{,}824 & 6.1\% \\
 & 4 & 14{,}555 & 14{,}555 & 1.0000 & 13{,}824 & 13{,}824 & 4.0\% \\
 & 5 & 11{,}828 & 11{,}828 & 1.0000 & 20{,}736 & 20{,}736 & 8.8\% \\
\addlinespace[2pt]
case\_023 & 1 & 107{,}465 & 107{,}331 & 1.0012 & 1{,}033{,}228 & 1{,}033{,}228 & 10.2\% \\
 & 2 & 58{,}256 & 58{,}234 & 1.0004 & 1{,}192{,}892 & 1{,}192{,}892 & 15.6\% \\
 & 3 & 40{,}897 & 40{,}897 & 1.0000 & 196{,}896 & 196{,}896 & 31.2\% \\
 & 4 & 30{,}755 & 30{,}755 & 1.0000 & 281{,}280 & 281{,}280 & 25.6\% \\
 & 5 & 24{,}947 & 24{,}947 & 1.0000 & 270{,}912 & 270{,}912 & 24.5\% \\
\addlinespace[2pt]
case\_024 & 1 & 2{,}555{,}343 & 2{,}555{,}343 & 1.0000 & 0 & 0 & 0.0\% \\
 & 2 & 2{,}334{,}112 & 2{,}334{,}112 & 1.0000 & 230{,}200 & 230{,}200 & 36.8\% \\
 & 3 & 1{,}062{,}198 & 1{,}062{,}198 & 1.0000 & 6{,}040 & 6{,}040 & 0.1\% \\
 & 4 & 846{,}529 & 846{,}529 & 1.0000 & 7{,}652 & 7{,}652 & 0.1\% \\
 & 5 & 846{,}529 & 846{,}529 & 1.0000 & 7{,}652 & 7{,}652 & 0.1\% \\
\addlinespace[2pt]
case\_025 & 1 & 4{,}863{,}026 & 4{,}863{,}026 & 1.0000 & 0 & 0 & 0.0\% \\
 & 2 & 2{,}431{,}840 & 2{,}431{,}840 & 1.0000 & 27{,}648 & 27{,}648 & 0.0\% \\
 & 3 & 1{,}622{,}164 & 1{,}622{,}164 & 1.0000 & 55{,}296 & 55{,}296 & 0.0\% \\
 & 4 & 1{,}217{,}444 & 1{,}217{,}444 & 1.0000 & 82{,}944 & 82{,}944 & 0.0\% \\
 & 5 & 975{,}282 & 975{,}282 & 1.0000 & 110{,}592 & 110{,}592 & 0.0\% \\
\addlinespace[2pt]
case\_026 & 1 & 96{,}504 & 96{,}504 & 1.0000 & 0 & 0 & 0.0\% \\
 & 2 & 47{,}565 & 47{,}565 & 1.0000 & 1{,}296 & 1{,}296 & 0.3\% \\
 & 3 & 41{,}536 & 41{,}536 & 1.0000 & 95{,}664 & 95{,}664 & 16.4\% \\
 & 4 & 41{,}536 & 41{,}536 & 1.0000 & 95{,}664 & 95{,}664 & 16.4\% \\
 & 5 & 38{,}400 & 36{,}058 & 1.0650 & 1{,}162{,}436 & 1{,}162{,}436 & 30.1\% \\
\addlinespace[2pt]
case\_027 & 1 & 2{,}363{,}952 & 2{,}359{,}515 & 1.0019 & 1{,}400{,}832 & 1{,}400{,}832 & 26.9\% \\
 & 2 & 1{,}120{,}962 & 1{,}120{,}962 & 1.0000 & 1{,}744{,}896 & 1{,}744{,}896 & 8.9\% \\
 & 3 & 837{,}049 & 837{,}049 & 1.0000 & 3{,}686{,}400 & 3{,}686{,}400 & 43.0\% \\
 & 4 & 559{,}196 & 559{,}196 & 1.0000 & 3{,}317{,}760 & 3{,}317{,}760 & 0.0\% \\
 & 5 & 558{,}753 & 558{,}753 & 1.0000 & 5{,}529{,}600 & 5{,}529{,}600 & 4.4\% \\
\addlinespace[2pt]
case\_028 & 1 & 40{,}787{,}625 & 40{,}787{,}625 & 1.0000 & 32{,}169{,}984 & 32{,}169{,}984 & 0.0\% \\
 & 2 & 20{,}475{,}783 & 20{,}471{,}568 & 1.0002 & 25{,}833{,}600 & 25{,}833{,}600 & 11.9\% \\
 & 3 & 13{,}590{,}237 & 13{,}590{,}237 & 1.0000 & 23{,}420{,}160 & 23{,}420{,}160 & 0.0\% \\
 & 4 & 10{,}268{,}731 & 10{,}247{,}047 & 1.0021 & 36{,}349{,}824 & 36{,}349{,}824 & 18.7\% \\
 & 5 & 8{,}169{,}383 & 8{,}146{,}739 & 1.0028 & 33{,}460{,}224 & 33{,}460{,}224 & 34.3\% \\
\addlinespace[2pt]
case\_029 & 1 & 37{,}555 & 37{,}555 & 1.0000 & 0 & 0 & 0.0\% \\
 & 2 & 19{,}740 & 19{,}740 & 1.0000 & 0 & 0 & 0.0\% \\
 & 3 & 12{,}743 & 12{,}743 & 1.0000 & 0 & 0 & 0.0\% \\
 & 4 & 11{,}028 & 11{,}028 & 1.0000 & 0 & 0 & 0.0\% \\
 & 5 & 9{,}383 & 9{,}383 & 1.0000 & 0 & 0 & 0.0\% \\
\addlinespace[2pt]
case\_030 & 1 & 2{,}808{,}034 & 2{,}807{,}891 & 1.0001 & 5{,}519{,}360 & 5{,}519{,}360 & 26.0\% \\
 & 2 & 1{,}390{,}756 & 1{,}390{,}756 & 1.0000 & 2{,}045{,}952 & 2{,}045{,}952 & 43.2\% \\
 & 3 & 937{,}837 & 936{,}307 & 1.0016 & 8{,}459{,}264 & 8{,}459{,}264 & 54.4\% \\
 & 4 & 696{,}635 & 696{,}481 & 1.0002 & 4{,}763{,}648 & 4{,}763{,}648 & 62.2\% \\
 & 5 & 560{,}819 & 560{,}766 & 1.0001 & 10{,}660{,}864 & 10{,}660{,}864 & 55.2\% \\
\addlinespace[2pt]
case\_031 & 1 & 1{,}414{,}369 & 1{,}414{,}159 & 1.0001 & 3{,}491{,}328 & 3{,}491{,}328 & 26.1\% \\
 & 2 & 734{,}985 & 734{,}098 & 1.0012 & 3{,}924{,}480 & 4{,}815{,}360 & 49.0\% \\
 & 3 & 486{,}170 & 485{,}410 & 1.0016 & 4{,}591{,}104 & 4{,}591{,}104 & 38.2\% \\
 & 4 & 390{,}627 & 390{,}627 & 1.0000 & 5{,}976{,}868 & 5{,}976{,}868 & 63.5\% \\
 & 5 & 307{,}985 & 307{,}871 & 1.0004 & 7{,}549{,}440 & 8{,}059{,}392 & 43.7\% \\
\addlinespace[2pt]
case\_032 & 1 & 34{,}557 & 34{,}557 & 1.0000 & 0 & 0 & 0.0\% \\
 & 2 & 19{,}737 & 19{,}737 & 1.0000 & 6{,}216 & 6{,}216 & 2.2\% \\
 & 3 & 12{,}854 & 12{,}854 & 1.0000 & 12{,}360 & 12{,}360 & 4.2\% \\
 & 4 & 12{,}854 & 12{,}854 & 1.0000 & 12{,}360 & 12{,}360 & 4.2\% \\
 & 5 & 10{,}415 & 10{,}378 & 1.0036 & 12{,}360 & 12{,}360 & 4.2\% \\
\addlinespace[2pt]
case\_033 & 1 & 1{,}774{,}196 & 1{,}774{,}080 & 1.0001 & 4{,}592{,}640 & 4{,}592{,}640 & 9.6\% \\
 & 2 & 887{,}381 & 886{,}959 & 1.0005 & 6{,}501{,}408 & 6{,}885{,}408 & 31.3\% \\
 & 3 & 602{,}654 & 602{,}341 & 1.0005 & 6{,}600{,}768 & 6{,}600{,}768 & 33.2\% \\
 & 4 & 459{,}481 & 457{,}965 & 1.0033 & 7{,}844{,}448 & 7{,}718{,}496 & 34.7\% \\
 & 5 & 376{,}887 & 372{,}694 & 1.0113 & 7{,}767{,}680 & 7{,}570{,}560 & 34.8\% \\
\addlinespace[2pt]
case\_034 & 1 & 825{,}167 & 825{,}167 & 1.0000 & 163{,}844 & 163{,}844 & 6.7\% \\
 & 2 & 412{,}835 & 412{,}796 & 1.0001 & 1{,}082{,}880 & 1{,}082{,}880 & 32.1\% \\
 & 3 & 271{,}887 & 271{,}520 & 1.0014 & 2{,}453{,}504 & 2{,}453{,}504 & 50.8\% \\
 & 4 & 207{,}388 & 206{,}225 & 1.0056 & 3{,}150{,}336 & 3{,}150{,}336 & 52.2\% \\
 & 5 & 170{,}400 & 167{,}446 & 1.0176 & 3{,}487{,}744 & 3{,}487{,}744 & 51.5\% \\
\addlinespace[2pt]
case\_035 & 1 & 156{,}657 & 156{,}657 & 1.0000 & 579{,}992 & 579{,}992 & 13.2\% \\
 & 2 & 95{,}406 & 95{,}406 & 1.0000 & 0 & 0 & 0.0\% \\
 & 3 & 95{,}406 & 95{,}406 & 1.0000 & 0 & 0 & 0.0\% \\
 & 4 & 95{,}406 & 95{,}406 & 1.0000 & 0 & 0 & 0.0\% \\
 & 5 & 95{,}406 & 95{,}406 & 1.0000 & 0 & 0 & 0.0\% \\
\addlinespace[2pt]
case\_036 & 1 & 931{,}510 & 931{,}510 & 1.0000 & 0 & 0 & 0.0\% \\
 & 2 & 466{,}068 & 466{,}068 & 1.0000 & 1{,}152 & 1{,}152 & 0.0\% \\
 & 3 & 311{,}345 & 311{,}345 & 1.0000 & 2{,}304 & 2{,}304 & 0.0\% \\
 & 4 & 233{,}584 & 233{,}584 & 1.0000 & 3{,}456 & 3{,}456 & 0.0\% \\
 & 5 & 187{,}182 & 187{,}182 & 1.0000 & 4{,}608 & 4{,}608 & 0.0\% \\
\addlinespace[2pt]
case\_037 & 1 & 210{,}001 & 210{,}001 & 1.0000 & 0 & 0 & 0.0\% \\
 & 2 & 105{,}374 & 105{,}374 & 1.0000 & 204{,}800 & 204{,}800 & 0.0\% \\
 & 3 & 74{,}038 & 74{,}038 & 1.0000 & 1{,}064{,}960 & 1{,}064{,}960 & 59.0\% \\
 & 4 & 53{,}320 & 53{,}320 & 1.0000 & 614{,}400 & 614{,}400 & 0.0\% \\
 & 5 & 43{,}800 & 43{,}013 & 1.0183 & 1{,}925{,}120 & 1{,}925{,}120 & 64.2\% \\
\addlinespace[2pt]
case\_038 & 1 & 1{,}364{,}555 & 1{,}363{,}312 & 1.0009 & 5{,}627{,}904 & 5{,}627{,}904 & 17.4\% \\
 & 2 & 678{,}569 & 678{,}530 & 1.0001 & 1{,}992{,}192 & 1{,}992{,}192 & 39.2\% \\
 & 3 & 456{,}517 & 454{,}899 & 1.0036 & 5{,}171{,}712 & 5{,}171{,}712 & 58.3\% \\
 & 4 & 346{,}741 & 344{,}588 & 1.0062 & 8{,}735{,}232 & 6{,}923{,}776 & 50.7\% \\
 & 5 & 285{,}550 & 282{,}185 & 1.0119 & 8{,}265{,}216 & 7{,}371{,}264 & 68.0\% \\
\addlinespace[2pt]
case\_039 & 1 & 1{,}435{,}262 & 1{,}435{,}262 & 1.0000 & 7{,}929{,}856 & 7{,}929{,}856 & 0.0\% \\
 & 2 & 708{,}561 & 705{,}615 & 1.0042 & 11{,}067{,}392 & 11{,}067{,}392 & 37.0\% \\
 & 3 & 485{,}565 & 483{,}226 & 1.0048 & 14{,}323{,}712 & 14{,}323{,}712 & 38.6\% \\
 & 4 & 395{,}323 & 368{,}055 & 1.0741 & 17{,}604{,}608 & 17{,}604{,}608 & 47.8\% \\
 & 5 & 395{,}323 & 340{,}476 & 1.1611 & 17{,}604{,}608 & 20{,}480{,}000 & 37.1\% \\
\addlinespace[2pt]
case\_040 & 1 & 253{,}454 & 253{,}454 & 1.0000 & 0 & 0 & 0.0\% \\
 & 2 & 120{,}986 & 120{,}986 & 1.0000 & 12{,}288 & 12{,}288 & 1.3\% \\
 & 3 & 100{,}177 & 100{,}177 & 1.0000 & 65{,}090 & 65{,}090 & 1.1\% \\
 & 4 & 99{,}511 & 99{,}511 & 1.0000 & 265{,}034 & 265{,}034 & 4.5\% \\
 & 5 & 95{,}285 & 95{,}283 & 1.0000 & 388{,}940 & 388{,}940 & 10.2\% \\
\addlinespace[2pt]
case\_041 & 1 & 6{,}003{,}957 & 6{,}001{,}575 & 1.0004 & 30{,}898{,}176 & 30{,}898{,}176 & 10.2\% \\
 & 2 & 2{,}989{,}434 & 2{,}986{,}558 & 1.0010 & 26{,}674{,}176 & 26{,}674{,}176 & 22.8\% \\
 & 3 & 2{,}006{,}212 & 2{,}001{,}038 & 1.0026 & 33{,}024{,}000 & 33{,}024{,}000 & 29.2\% \\
 & 4 & 1{,}497{,}746 & 1{,}494{,}017 & 1.0025 & 34{,}397{,}184 & 34{,}397{,}184 & 24.1\% \\
 & 5 & 1{,}203{,}317 & 1{,}199{,}853 & 1.0029 & 36{,}360{,}192 & 31{,}334{,}400 & 38.5\% \\
\addlinespace[2pt]
case\_042 & 1 & 1{,}686{,}740 & 1{,}675{,}633 & 1.0066 & 6{,}629{,}376 & 6{,}629{,}376 & 31.9\% \\
 & 2 & 845{,}035 & 844{,}570 & 1.0006 & 6{,}009{,}120 & 6{,}540{,}576 & 37.1\% \\
 & 3 & 584{,}180 & 581{,}304 & 1.0049 & 8{,}561{,}952 & 8{,}396{,}064 & 34.5\% \\
 & 4 & 438{,}568 & 437{,}100 & 1.0034 & 8{,}294{,}688 & 8{,}294{,}688 & 36.8\% \\
 & 5 & 355{,}113 & 354{,}356 & 1.0021 & 10{,}383{,}360 & 10{,}383{,}360 & 43.3\% \\
\addlinespace[2pt]
case\_043 & 1 & 674{,}823 & 674{,}823 & 1.0000 & 16{,}384 & 16{,}384 & 1.1\% \\
 & 2 & 338{,}088 & 338{,}088 & 1.0000 & 460{,}800 & 460{,}800 & 21.9\% \\
 & 3 & 245{,}028 & 245{,}028 & 1.0000 & 352{,}450 & 352{,}450 & 10.8\% \\
 & 4 & 245{,}028 & 245{,}028 & 1.0000 & 352{,}450 & 352{,}450 & 10.8\% \\
 & 5 & 233{,}701 & 223{,}774 & 1.0444 & 7{,}587{,}124 & 7{,}587{,}124 & 21.4\% \\
\addlinespace[2pt]
case\_044 & 1 & 74{,}811 & 74{,}811 & 1.0000 & 0 & 0 & 0.0\% \\
 & 2 & 43{,}931 & 41{,}665 & 1.0544 & 930{,}400 & 930{,}400 & 14.8\% \\
 & 3 & 43{,}931 & 41{,}665 & 1.0544 & 930{,}400 & 930{,}400 & 14.8\% \\
 & 4 & 43{,}931 & 41{,}665 & 1.0544 & 930{,}400 & 930{,}400 & 14.8\% \\
 & 5 & 43{,}931 & 41{,}665 & 1.0544 & 930{,}400 & 930{,}400 & 14.8\% \\
\addlinespace[2pt]
case\_045 & 1 & 125{,}614 & 125{,}614 & 1.0000 & 0 & 0 & 0.0\% \\
 & 2 & 63{,}059 & 63{,}059 & 1.0000 & 0 & 0 & 0.0\% \\
 & 3 & 42{,}601 & 42{,}601 & 1.0000 & 0 & 0 & 0.0\% \\
 & 4 & 31{,}878 & 31{,}878 & 1.0000 & 0 & 0 & 0.0\% \\
 & 5 & 26{,}164 & 26{,}164 & 1.0000 & 0 & 0 & 0.0\% \\
\addlinespace[2pt]
case\_046 & 1 & 210{,}387 & 210{,}387 & 1.0000 & 0 & 0 & 0.0\% \\
 & 2 & 106{,}042 & 106{,}042 & 1.0000 & 930{,}400 & 930{,}400 & 0.0\% \\
 & 3 & 80{,}028 & 80{,}028 & 1.0000 & 1{,}860{,}800 & 1{,}860{,}800 & 63.8\% \\
 & 4 & 78{,}210 & 75{,}972 & 1.0295 & 2{,}791{,}200 & 3{,}181{,}344 & 30.9\% \\
 & 5 & 78{,}210 & 75{,}972 & 1.0295 & 2{,}791{,}200 & 3{,}181{,}344 & 30.9\% \\
\addlinespace[2pt]
case\_047 & 1 & 291{,}968 & 291{,}968 & 1.0000 & 0 & 0 & 0.0\% \\
 & 2 & 177{,}730 & 177{,}702 & 1.0002 & 2{,}253{,}222 & 2{,}253{,}222 & 8.1\% \\
 & 3 & 138{,}079 & 137{,}526 & 1.0040 & 3{,}288{,}588 & 3{,}288{,}588 & 13.5\% \\
 & 4 & 128{,}508 & 126{,}690 & 1.0143 & 4{,}357{,}842 & 4{,}357{,}842 & 18.7\% \\
 & 5 & 128{,}508 & 123{,}689 & 1.0390 & 4{,}357{,}842 & 5{,}448{,}792 & 26.1\% \\
\addlinespace[2pt]
case\_048 & 1 & 190{,}989 & 190{,}989 & 1.0000 & 0 & 0 & 0.0\% \\
 & 2 & 129{,}426 & 129{,}426 & 1.0000 & 314{,}182 & 314{,}182 & 19.0\% \\
 & 3 & 113{,}494 & 113{,}494 & 1.0000 & 456{,}716 & 456{,}716 & 16.7\% \\
 & 4 & 111{,}484 & 111{,}311 & 1.0016 & 1{,}039{,}728 & 1{,}039{,}728 & 9.8\% \\
 & 5 & 111{,}484 & 111{,}311 & 1.0016 & 1{,}039{,}728 & 1{,}039{,}728 & 9.8\% \\
\addlinespace[2pt]
case\_049 & 1 & 175{,}821 & 175{,}821 & 1.0000 & 11{,}280 & 11{,}280 & 0.0\% \\
 & 2 & 126{,}134 & 126{,}132 & 1.0000 & 1{,}316{,}902 & 1{,}316{,}902 & 2.0\% \\
 & 3 & 99{,}627 & 99{,}620 & 1.0001 & 1{,}901{,}512 & 1{,}901{,}512 & 7.2\% \\
 & 4 & 89{,}606 & 89{,}332 & 1.0031 & 2{,}249{,}534 & 2{,}249{,}534 & 11.9\% \\
 & 5 & 89{,}606 & 89{,}332 & 1.0031 & 2{,}249{,}534 & 2{,}249{,}534 & 11.9\% \\
\addlinespace[2pt]
case\_050 & 1 & 142{,}868 & 142{,}181 & 1.0048 & 681{,}388 & 681{,}388 & 15.1\% \\
 & 2 & 84{,}490 & 84{,}225 & 1.0031 & 854{,}586 & 854{,}586 & 22.5\% \\
 & 3 & 72{,}131 & 68{,}699 & 1.0500 & 1{,}528{,}238 & 1{,}528{,}238 & 19.8\% \\
 & 4 & 64{,}647 & 60{,}156 & 1.0747 & 1{,}809{,}716 & 1{,}809{,}716 & 27.9\% \\
 & 5 & 64{,}647 & 60{,}156 & 1.0747 & 1{,}809{,}716 & 1{,}809{,}716 & 27.9\% \\
\addlinespace[2pt]
case\_051 & 1 & 607{,}628 & 607{,}628 & 1.0000 & 0 & 0 & 0.0\% \\
 & 2 & 547{,}789 & 547{,}789 & 1.0000 & 229{,}584 & 229{,}584 & 36.8\% \\
 & 3 & 252{,}851 & 252{,}851 & 1.0000 & 1{,}420 & 1{,}420 & 0.0\% \\
 & 4 & 202{,}281 & 202{,}281 & 1.0000 & 1{,}724 & 1{,}724 & 0.0\% \\
 & 5 & 201{,}161 & 201{,}161 & 1.0000 & 1{,}808 & 1{,}808 & 0.0\% \\
\addlinespace[2pt]
case\_052 & 1 & 179{,}248 & 179{,}248 & 1.0000 & 0 & 0 & 0.0\% \\
 & 2 & 92{,}135 & 91{,}913 & 1.0024 & 1{,}938{,}984 & 1{,}938{,}984 & 10.1\% \\
 & 3 & 68{,}329 & 68{,}129 & 1.0029 & 1{,}409{,}612 & 2{,}403{,}994 & 15.3\% \\
 & 4 & 61{,}312 & 60{,}923 & 1.0064 & 259{,}584 & 2{,}174{,}152 & 16.1\% \\
 & 5 & 50{,}551 & 50{,}547 & 1.0001 & 328{,}704 & 328{,}704 & 19.6\% \\
\addlinespace[2pt]
case\_053 & 1 & 1{,}293{,}873 & 1{,}293{,}873 & 1.0000 & 478{,}218 & 478{,}218 & 0.4\% \\
 & 2 & 698{,}368 & 694{,}294 & 1.0059 & 8{,}085{,}394 & 8{,}085{,}394 & 14.2\% \\
 & 3 & 506{,}815 & 506{,}815 & 1.0000 & 5{,}882{,}944 & 5{,}882{,}944 & 58.1\% \\
 & 4 & 477{,}534 & 460{,}610 & 1.0367 & 6{,}416{,}288 & 17{,}224{,}714 & 20.9\% \\
 & 5 & 443{,}726 & 421{,}386 & 1.0530 & 15{,}484{,}744 & 15{,}484{,}744 & 20.2\% \\
\addlinespace[2pt]
case\_054 & 1 & 814{,}418 & 814{,}418 & 1.0000 & 28{,}160 & 28{,}160 & 0.0\% \\
 & 2 & 472{,}156 & 472{,}156 & 1.0000 & 6{,}564{,}694 & 6{,}564{,}694 & 2.4\% \\
 & 3 & 347{,}229 & 346{,}008 & 1.0035 & 9{,}468{,}104 & 9{,}468{,}104 & 9.0\% \\
 & 4 & 317{,}609 & 316{,}029 & 1.0050 & 7{,}167{,}478 & 9{,}639{,}186 & 14.9\% \\
 & 5 & 287{,}243 & 278{,}185 & 1.0326 & 11{,}160{,}124 & 11{,}160{,}124 & 20.9\% \\
\addlinespace[2pt]
case\_055 & 1 & 111{,}355 & 110{,}071 & 1.0117 & 1{,}509{,}616 & 1{,}509{,}616 & 11.7\% \\
 & 2 & 56{,}668 & 56{,}563 & 1.0019 & 57{,}856 & 57{,}856 & 8.5\% \\
 & 3 & 38{,}406 & 38{,}406 & 1.0000 & 115{,}712 & 115{,}712 & 15.6\% \\
 & 4 & 31{,}367 & 31{,}367 & 1.0000 & 115{,}712 & 115{,}712 & 15.6\% \\
 & 5 & 26{,}322 & 26{,}322 & 1.0000 & 198{,}656 & 198{,}656 & 24.1\% \\
\addlinespace[2pt]
case\_056 & 1 & 252{,}916 & 252{,}916 & 1.0000 & 145{,}096 & 145{,}096 & 13.8\% \\
 & 2 & 172{,}161 & 171{,}946 & 1.0013 & 1{,}922{,}368 & 1{,}922{,}368 & 8.4\% \\
 & 3 & 129{,}190 & 128{,}913 & 1.0021 & 2{,}780{,}198 & 2{,}780{,}198 & 11.4\% \\
 & 4 & 124{,}871 & 123{,}907 & 1.0078 & 3{,}387{,}698 & 3{,}387{,}698 & 17.6\% \\
 & 5 & 124{,}871 & 121{,}373 & 1.0288 & 3{,}387{,}698 & 4{,}645{,}072 & 26.5\% \\
\addlinespace[2pt]
case\_057 & 1 & 58{,}977 & 58{,}977 & 1.0000 & 20{,}484 & 20{,}484 & 5.7\% \\
 & 2 & 33{,}603 & 33{,}603 & 1.0000 & 288 & 288 & 0.1\% \\
 & 3 & 25{,}250 & 25{,}250 & 1.0000 & 288 & 288 & 0.1\% \\
 & 4 & 20{,}957 & 20{,}957 & 1.0000 & 288 & 288 & 0.1\% \\
 & 5 & 16{,}538 & 16{,}538 & 1.0000 & 27{,}936 & 27{,}936 & 7.7\% \\
\addlinespace[2pt]
case\_058 & 1 & 4{,}539{,}717 & 4{,}539{,}717 & 1.0000 & 28{,}696{,}576 & 28{,}696{,}576 & 0.0\% \\
 & 2 & 2{,}271{,}316 & 2{,}270{,}206 & 1.0005 & 33{,}431{,}552 & 33{,}431{,}552 & 9.9\% \\
 & 3 & 1{,}554{,}831 & 1{,}554{,}831 & 1.0000 & 30{,}400{,}512 & 30{,}400{,}512 & 0.0\% \\
 & 4 & 1{,}154{,}519 & 1{,}142{,}629 & 1.0104 & 42{,}622{,}976 & 42{,}622{,}976 & 26.9\% \\
 & 5 & 1{,}004{,}819 & 1{,}004{,}819 & 1.0000 & 32{,}104{,}448 & 32{,}104{,}448 & 0.0\% \\
\addlinespace[2pt]
case\_059 & 1 & 829{,}473 & 815{,}938 & 1.0166 & 1{,}691{,}648 & 1{,}691{,}648 & 61.0\% \\
 & 2 & 369{,}318 & 369{,}318 & 1.0000 & 360{,}448 & 360{,}448 & 0.0\% \\
 & 3 & 266{,}217 & 262{,}333 & 1.0148 & 2{,}998{,}272 & 2{,}998{,}272 & 68.3\% \\
 & 4 & 185{,}216 & 185{,}216 & 1.0000 & 1{,}081{,}344 & 1{,}081{,}344 & 0.0\% \\
 & 5 & 150{,}782 & 150{,}736 & 1.0003 & 4{,}325{,}376 & 4{,}325{,}376 & 73.8\% \\
\addlinespace[2pt]
case\_060 & 1 & 883{,}206 & 882{,}920 & 1.0003 & 1{,}539{,}108 & 1{,}539{,}108 & 30.6\% \\
 & 2 & 464{,}279 & 464{,}106 & 1.0004 & 1{,}705{,}984 & 2{,}004{,}992 & 33.0\% \\
 & 3 & 334{,}694 & 334{,}668 & 1.0001 & 2{,}796{,}672 & 2{,}796{,}672 & 42.0\% \\
 & 4 & 282{,}484 & 282{,}423 & 1.0002 & 3{,}623{,}168 & 3{,}623{,}168 & 42.6\% \\
 & 5 & 237{,}234 & 237{,}114 & 1.0005 & 4{,}286{,}848 & 4{,}286{,}848 & 44.9\% \\
\addlinespace[2pt]
case\_061 & 1 & 77{,}191 & 77{,}191 & 1.0000 & 0 & 0 & 0.0\% \\
 & 2 & 47{,}568 & 47{,}372 & 1.0041 & 671{,}274 & 671{,}274 & 11.8\% \\
 & 3 & 28{,}409 & 28{,}368 & 1.0014 & 43{,}008 & 43{,}008 & 7.0\% \\
 & 4 & 22{,}111 & 22{,}111 & 1.0000 & 62{,}208 & 62{,}208 & 9.2\% \\
 & 5 & 18{,}783 & 18{,}751 & 1.0017 & 103{,}936 & 103{,}936 & 14.4\% \\
\addlinespace[2pt]
case\_062 & 1 & 2{,}854{,}773 & 2{,}854{,}773 & 1.0000 & 0 & 0 & 0.0\% \\
 & 2 & 1{,}718{,}716 & 1{,}717{,}254 & 1.0009 & 62{,}917{,}632 & 62{,}917{,}632 & 0.7\% \\
 & 3 & 1{,}056{,}247 & 1{,}056{,}247 & 1.0000 & 8{,}732{,}160 & 8{,}732{,}160 & 0.0\% \\
 & 4 & 1{,}056{,}247 & 1{,}056{,}247 & 1.0000 & 8{,}732{,}160 & 8{,}732{,}160 & 0.0\% \\
 & 5 & 1{,}056{,}247 & 1{,}056{,}247 & 1.0000 & 8{,}732{,}160 & 8{,}732{,}160 & 0.0\% \\
\addlinespace[2pt]
case\_063 & 1 & 1{,}070{,}361 & 1{,}070{,}361 & 1.0000 & 0 & 0 & 0.0\% \\
 & 2 & 618{,}185 & 625{,}947 & 0.9876 & 20{,}963{,}328 & 20{,}949{,}504 & 0.9\% \\
 & 3 & 397{,}260 & 397{,}260 & 1.0000 & 3{,}081{,}216 & 3{,}081{,}216 & 0.0\% \\
 & 4 & 301{,}177 & 301{,}177 & 1.0000 & 2{,}672{,}640 & 2{,}672{,}640 & 0.0\% \\
 & 5 & 244{,}345 & 244{,}345 & 1.0000 & 2{,}939{,}904 & 2{,}939{,}904 & 0.0\% \\
\addlinespace[2pt]
case\_064 & 1 & 17{,}034 & 17{,}034 & 1.0000 & 0 & 0 & 0.0\% \\
 & 2 & 13{,}414 & 13{,}414 & 1.0000 & 77{,}222 & 77{,}222 & 9.2\% \\
 & 3 & 12{,}280 & 12{,}280 & 1.0000 & 94{,}666 & 94{,}666 & 21.4\% \\
 & 4 & 11{,}941 & 11{,}942 & 0.9999 & 143{,}022 & 143{,}022 & 24.9\% \\
 & 5 & 11{,}406 & 11{,}391 & 1.0013 & 142{,}706 & 142{,}706 & 27.2\% \\
\addlinespace[2pt]
case\_065 & 1 & 47{,}392 & 47{,}392 & 1.0000 & 0 & 0 & 0.0\% \\
 & 2 & 25{,}997 & 25{,}997 & 1.0000 & 13{,}968 & 13{,}968 & 4.5\% \\
 & 3 & 18{,}001 & 18{,}001 & 1.0000 & 27{,}936 & 27{,}936 & 5.7\% \\
 & 4 & 14{,}178 & 14{,}178 & 1.0000 & 41{,}904 & 41{,}904 & 8.2\% \\
 & 5 & 12{,}135 & 12{,}135 & 1.0000 & 55{,}584 & 55{,}584 & 15.6\% \\
\addlinespace[2pt]
case\_066 & 1 & 336{,}224 & 336{,}224 & 1.0000 & 123{,}762 & 123{,}762 & 0.0\% \\
 & 2 & 183{,}389 & 183{,}351 & 1.0002 & 1{,}932{,}976 & 1{,}932{,}976 & 13.9\% \\
 & 3 & 128{,}156 & 127{,}738 & 1.0033 & 2{,}549{,}320 & 2{,}549{,}320 & 26.3\% \\
 & 4 & 110{,}097 & 109{,}279 & 1.0075 & 3{,}362{,}510 & 3{,}362{,}510 & 21.1\% \\
 & 5 & 109{,}083 & 102{,}063 & 1.0688 & 3{,}836{,}102 & 4{,}044{,}374 & 26.6\% \\
\addlinespace[2pt]
case\_067 & 1 & 61{,}679{,}114 & 61{,}412{,}609 & 1.0043 & 51{,}113{,}984 & 194{,}226{,}176 & 92.3\% \\
 & 2 & 30{,}476{,}248 & 30{,}437{,}438 & 1.0013 & 56{,}984{,}192 & 56{,}984{,}192 & 13.5\% \\
 & 3 & 20{,}333{,}946 & 20{,}264{,}004 & 1.0035 & 17{,}584{,}896 & 17{,}584{,}896 & 42.8\% \\
 & 4 & 15{,}271{,}913 & 15{,}168{,}347 & 1.0068 & 16{,}931{,}968 & 16{,}931{,}968 & 59.0\% \\
 & 5 & 12{,}621{,}657 & 12{,}587{,}135 & 1.0027 & 17{,}327{,}616 & 17{,}327{,}616 & 17.7\% \\
\addlinespace[2pt]
case\_068 & 1 & 342{,}841 & 342{,}271 & 1.0017 & 518{,}916 & 518{,}916 & 4.0\% \\
 & 2 & 202{,}721 & 202{,}474 & 1.0012 & 2{,}254{,}270 & 2{,}254{,}270 & 13.9\% \\
 & 3 & 157{,}301 & 154{,}308 & 1.0194 & 3{,}594{,}672 & 3{,}594{,}672 & 17.5\% \\
 & 4 & 143{,}171 & 132{,}760 & 1.0784 & 4{,}306{,}378 & 3{,}365{,}042 & 22.9\% \\
 & 5 & 143{,}025 & 132{,}760 & 1.0773 & 3{,}353{,}518 & 3{,}365{,}042 & 22.9\% \\
\addlinespace[2pt]
case\_069 & 1 & 21{,}794 & 21{,}794 & 1.0000 & 0 & 0 & 0.0\% \\
 & 2 & 13{,}965 & 13{,}887 & 1.0056 & 404{,}806 & 404{,}806 & 9.9\% \\
 & 3 & 13{,}965 & 13{,}887 & 1.0056 & 404{,}806 & 404{,}806 & 9.9\% \\
 & 4 & 13{,}965 & 13{,}887 & 1.0056 & 404{,}806 & 404{,}806 & 9.9\% \\
 & 5 & 13{,}831 & 13{,}471 & 1.0267 & 377{,}328 & 377{,}328 & 28.6\% \\
\addlinespace[2pt]
case\_070 & 1 & 112{,}310 & 112{,}310 & 1.0000 & 0 & 0 & 0.0\% \\
 & 2 & 56{,}421 & 56{,}410 & 1.0002 & 129{,}536 & 129{,}536 & 19.1\% \\
 & 3 & 40{,}990 & 40{,}964 & 1.0006 & 259{,}072 & 259{,}072 & 9.7\% \\
 & 4 & 33{,}697 & 33{,}636 & 1.0018 & 360{,}960 & 360{,}960 & 20.7\% \\
 & 5 & 30{,}598 & 30{,}501 & 1.0032 & 462{,}848 & 1{,}257{,}704 & 32.2\% \\
\addlinespace[2pt]
case\_071 & 1 & 16{,}351 & 16{,}351 & 1.0000 & 0 & 0 & 0.0\% \\
 & 2 & 11{,}580 & 11{,}571 & 1.0008 & 269{,}862 & 269{,}862 & 8.2\% \\
 & 3 & 11{,}188 & 11{,}070 & 1.0107 & 400{,}138 & 400{,}138 & 22.6\% \\
 & 4 & 11{,}188 & 11{,}070 & 1.0107 & 400{,}138 & 400{,}138 & 22.6\% \\
 & 5 & 11{,}188 & 11{,}070 & 1.0107 & 400{,}138 & 400{,}138 & 22.6\% \\
\addlinespace[2pt]
case\_072 & 1 & 23{,}897{,}828 & 23{,}890{,}756 & 1.0003 & 108{,}613{,}680 & 108{,}613{,}680 & 2.6\% \\
 & 2 & 11{,}899{,}903 & 11{,}876{,}654 & 1.0020 & 91{,}122{,}464 & 91{,}122{,}464 & 18.2\% \\
 & 3 & 8{,}008{,}124 & 7{,}980{,}888 & 1.0034 & 113{,}388{,}320 & 113{,}388{,}320 & 12.2\% \\
 & 4 & 5{,}993{,}260 & 5{,}988{,}828 & 1.0007 & 122{,}669{,}856 & 122{,}669{,}856 & 10.7\% \\
 & 5 & 4{,}861{,}592 & 4{,}855{,}697 & 1.0012 & 113{,}066{,}672 & 113{,}066{,}672 & 10.6\% \\
\addlinespace[2pt]
case\_073 & 1 & 8{,}214{,}707 & 8{,}214{,}707 & 1.0000 & 188{,}416 & 188{,}416 & 0.0\% \\
 & 2 & 4{,}238{,}435 & 4{,}199{,}635 & 1.0092 & 3{,}721{,}600 & 3{,}721{,}600 & 48.9\% \\
 & 3 & 2{,}825{,}285 & 2{,}800{,}629 & 1.0088 & 7{,}443{,}200 & 7{,}443{,}200 & 36.2\% \\
 & 4 & 2{,}225{,}291 & 2{,}125{,}171 & 1.0471 & 11{,}164{,}800 & 11{,}164{,}800 & 74.2\% \\
 & 5 & 1{,}849{,}407 & 1{,}723{,}485 & 1.0731 & 14{,}886{,}400 & 25{,}637{,}888 & 86.5\% \\
\addlinespace[2pt]
case\_074 & 1 & 1{,}682{,}239 & 1{,}678{,}120 & 1.0025 & 6{,}070{,}280 & 6{,}070{,}280 & 16.7\% \\
 & 2 & 811{,}260 & 810{,}928 & 1.0004 & 7{,}262{,}208 & 7{,}262{,}208 & 23.9\% \\
 & 3 & 542{,}279 & 542{,}130 & 1.0003 & 7{,}397{,}376 & 7{,}397{,}376 & 24.4\% \\
 & 4 & 409{,}816 & 409{,}294 & 1.0013 & 8{,}658{,}948 & 8{,}658{,}948 & 36.3\% \\
 & 5 & 341{,}700 & 335{,}218 & 1.0193 & 8{,}777{,}728 & 8{,}777{,}728 & 36.4\% \\
\addlinespace[2pt]
case\_075 & 1 & 1{,}247{,}823 & 1{,}247{,}823 & 1.0000 & 0 & 0 & 0.0\% \\
 & 2 & 689{,}383 & 688{,}753 & 1.0009 & 3{,}963{,}590 & 3{,}963{,}590 & 15.7\% \\
 & 3 & 570{,}824 & 570{,}775 & 1.0001 & 5{,}602{,}700 & 5{,}602{,}700 & 28.6\% \\
 & 4 & 408{,}170 & 407{,}975 & 1.0005 & 6{,}136{,}978 & 6{,}136{,}978 & 31.4\% \\
 & 5 & 377{,}686 & 374{,}326 & 1.0090 & 8{,}326{,}392 & 8{,}326{,}392 & 36.0\% \\
\addlinespace[2pt]
case\_076 & 1 & 6{,}185{,}113 & 6{,}183{,}098 & 1.0003 & 28{,}280{,}832 & 28{,}280{,}832 & 13.4\% \\
 & 2 & 3{,}094{,}537 & 3{,}089{,}084 & 1.0018 & 30{,}394{,}368 & 30{,}394{,}368 & 36.2\% \\
 & 3 & 2{,}068{,}186 & 2{,}060{,}547 & 1.0037 & 31{,}798{,}272 & 31{,}798{,}272 & 46.7\% \\
 & 4 & 1{,}535{,}669 & 1{,}534{,}453 & 1.0008 & 26{,}690{,}560 & 26{,}690{,}560 & 44.0\% \\
 & 5 & 1{,}246{,}135 & 1{,}234{,}358 & 1.0095 & 35{,}584{,}000 & 37{,}436{,}416 & 50.4\% \\
\addlinespace[2pt]
case\_077 & 1 & 494{,}940 & 494{,}940 & 1.0000 & 0 & 0 & 0.0\% \\
 & 2 & 249{,}722 & 249{,}722 & 1.0000 & 57{,}344 & 57{,}344 & 2.6\% \\
 & 3 & 177{,}575 & 177{,}575 & 1.0000 & 114{,}688 & 114{,}688 & 2.5\% \\
 & 4 & 145{,}044 & 145{,}044 & 1.0000 & 172{,}032 & 172{,}032 & 5.0\% \\
 & 5 & 120{,}068 & 120{,}068 & 1.0000 & 172{,}032 & 172{,}032 & 0.0\% \\
\addlinespace[2pt]
case\_078 & 1 & 2{,}065{,}085 & 2{,}065{,}085 & 1.0000 & 0 & 0 & 0.0\% \\
 & 2 & 1{,}024{,}634 & 1{,}024{,}634 & 1.0000 & 2{,}148{,}480 & 2{,}148{,}480 & 0.0\% \\
 & 3 & 768{,}769 & 765{,}147 & 1.0047 & 4{,}849{,}920 & 4{,}849{,}920 & 60.6\% \\
 & 4 & 520{,}320 & 520{,}320 & 1.0000 & 6{,}445{,}440 & 6{,}445{,}440 & 0.0\% \\
 & 5 & 520{,}320 & 511{,}638 & 1.0170 & 6{,}445{,}440 & 9{,}146{,}880 & 73.5\% \\
\addlinespace[2pt]
case\_079 & 1 & 31{,}052{,}251 & 31{,}052{,}251 & 1.0000 & 32{,}448{,}512 & 32{,}448{,}512 & 0.0\% \\
 & 2 & 15{,}446{,}486 & 15{,}446{,}486 & 1.0000 & 31{,}780{,}864 & 31{,}780{,}864 & 0.0\% \\
 & 3 & 10{,}181{,}692 & 10{,}175{,}315 & 1.0006 & 39{,}960{,}576 & 39{,}960{,}576 & 13.7\% \\
 & 4 & 7{,}635{,}624 & 7{,}635{,}624 & 1.0000 & 30{,}445{,}568 & 30{,}445{,}568 & 0.0\% \\
 & 5 & 6{,}075{,}976 & 6{,}075{,}472 & 1.0001 & 47{,}677{,}440 & 47{,}677{,}440 & 24.4\% \\
\addlinespace[2pt]
case\_080 & 1 & 292{,}138 & 291{,}826 & 1.0011 & 548{,}864 & 548{,}864 & 16.9\% \\
 & 2 & 145{,}926 & 145{,}674 & 1.0017 & 2{,}191{,}360 & 2{,}191{,}360 & 43.6\% \\
 & 3 & 100{,}325 & 100{,}325 & 1.0000 & 3{,}047{,}424 & 3{,}047{,}424 & 0.0\% \\
 & 4 & 100{,}325 & 90{,}584 & 1.1075 & 3{,}047{,}424 & 4{,}190{,}208 & 47.9\% \\
 & 5 & 100{,}325 & 90{,}408 & 1.1097 & 3{,}047{,}424 & 4{,}571{,}136 & 34.4\% \\
\addlinespace[2pt]
case\_081 & 1 & 151{,}795 & 151{,}795 & 1.0000 & 0 & 0 & 0.0\% \\
 & 2 & 76{,}860 & 76{,}860 & 1.0000 & 0 & 0 & 0.0\% \\
 & 3 & 51{,}943 & 51{,}943 & 1.0000 & 0 & 0 & 0.0\% \\
 & 4 & 39{,}588 & 39{,}588 & 1.0000 & 0 & 0 & 0.0\% \\
 & 5 & 30{,}943 & 30{,}943 & 1.0000 & 0 & 0 & 0.0\% \\
\addlinespace[2pt]
case\_082 & 1 & 686{,}840 & 686{,}840 & 1.0000 & 361{,}610 & 361{,}610 & 0.0\% \\
 & 2 & 366{,}243 & 366{,}242 & 1.0000 & 2{,}233{,}350 & 2{,}233{,}350 & 7.0\% \\
 & 3 & 264{,}390 & 264{,}378 & 1.0000 & 3{,}467{,}596 & 3{,}467{,}596 & 10.8\% \\
 & 4 & 229{,}458 & 229{,}335 & 1.0005 & 4{,}577{,}264 & 4{,}577{,}264 & 26.6\% \\
 & 5 & 204{,}292 & 202{,}383 & 1.0094 & 6{,}027{,}608 & 6{,}233{,}688 & 38.6\% \\
\addlinespace[2pt]
case\_083 & 1 & 840{,}797 & 840{,}797 & 1.0000 & 131{,}072 & 131{,}072 & 0.0\% \\
 & 2 & 419{,}777 & 419{,}777 & 1.0000 & 930{,}400 & 930{,}400 & 44.0\% \\
 & 3 & 289{,}139 & 289{,}139 & 1.0000 & 1{,}860{,}800 & 1{,}860{,}800 & 30.6\% \\
 & 4 & 211{,}087 & 211{,}087 & 1.0000 & 2{,}791{,}200 & 2{,}791{,}200 & 70.2\% \\
 & 5 & 185{,}314 & 185{,}314 & 1.0000 & 3{,}721{,}600 & 3{,}721{,}600 & 19.0\% \\
\addlinespace[2pt]
case\_084 & 1 & 2{,}507{,}412 & 2{,}507{,}412 & 1.0000 & 0 & 0 & 0.0\% \\
 & 2 & 1{,}255{,}052 & 1{,}255{,}052 & 1.0000 & 0 & 0 & 0.0\% \\
 & 3 & 836{,}220 & 836{,}220 & 1.0000 & 0 & 0 & 0.0\% \\
 & 4 & 629{,}028 & 629{,}028 & 1.0000 & 0 & 0 & 0.0\% \\
 & 5 & 503{,}472 & 503{,}472 & 1.0000 & 0 & 0 & 0.0\% \\
\addlinespace[2pt]
case\_085 & 1 & 3{,}184{,}152 & 3{,}184{,}152 & 1.0000 & 345{,}888 & 345{,}888 & 0.0\% \\
 & 2 & 1{,}743{,}831 & 1{,}743{,}831 & 1.0000 & 13{,}942{,}986 & 13{,}942{,}986 & 6.3\% \\
 & 3 & 1{,}246{,}519 & 1{,}245{,}549 & 1.0008 & 21{,}276{,}524 & 21{,}276{,}524 & 15.7\% \\
 & 4 & 946{,}039 & 942{,}396 & 1.0039 & 21{,}333{,}874 & 21{,}333{,}874 & 15.6\% \\
 & 5 & 846{,}587 & 833{,}297 & 1.0159 & 25{,}456{,}536 & 25{,}335{,}448 & 30.8\% \\
\addlinespace[2pt]
case\_086 & 1 & 106{,}479 & 106{,}479 & 1.0000 & 72{,}390 & 72{,}390 & 0.0\% \\
 & 2 & 65{,}985 & 65{,}965 & 1.0003 & 1{,}576{,}580 & 1{,}576{,}580 & 3.4\% \\
 & 3 & 57{,}154 & 56{,}070 & 1.0193 & 2{,}169{,}644 & 2{,}169{,}644 & 17.6\% \\
 & 4 & 56{,}538 & 54{,}363 & 1.0400 & 2{,}561{,}616 & 2{,}561{,}616 & 21.7\% \\
 & 5 & 56{,}538 & 54{,}363 & 1.0400 & 2{,}561{,}616 & 2{,}561{,}616 & 21.7\% \\
\addlinespace[2pt]
case\_087 & 1 & 3{,}974{,}071 & 3{,}969{,}779 & 1.0011 & 7{,}902{,}720 & 7{,}902{,}720 & 18.3\% \\
 & 2 & 1{,}985{,}234 & 1{,}984{,}074 & 1.0006 & 9{,}999{,}360 & 9{,}999{,}360 & 26.8\% \\
 & 3 & 1{,}379{,}480 & 1{,}379{,}498 & 1.0000 & 9{,}183{,}744 & 10{,}082{,}304 & 23.3\% \\
 & 4 & 1{,}069{,}661 & 1{,}069{,}661 & 1.0000 & 11{,}672{,}064 & 11{,}672{,}064 & 23.8\% \\
 & 5 & 885{,}851 & 885{,}832 & 1.0000 & 8{,}663{,}040 & 8{,}663{,}040 & 30.4\% \\
\addlinespace[2pt]
case\_088 & 1 & 228{,}471 & 228{,}471 & 1.0000 & 325{,}640 & 325{,}640 & 1.4\% \\
 & 2 & 133{,}265 & 133{,}219 & 1.0003 & 1{,}295{,}018 & 1{,}295{,}018 & 15.6\% \\
 & 3 & 105{,}066 & 103{,}383 & 1.0163 & 2{,}175{,}676 & 2{,}978{,}740 & 23.7\% \\
 & 4 & 90{,}485 & 87{,}105 & 1.0388 & 2{,}472{,}754 & 2{,}472{,}754 & 22.4\% \\
 & 5 & 90{,}485 & 87{,}105 & 1.0388 & 2{,}472{,}754 & 2{,}472{,}754 & 22.4\% \\
\addlinespace[2pt]
case\_089 & 1 & 595{,}298 & 595{,}298 & 1.0000 & 0 & 0 & 0.0\% \\
 & 2 & 295{,}856 & 295{,}856 & 1.0000 & 238{,}720 & 238{,}720 & 0.0\% \\
 & 3 & 208{,}395 & 208{,}395 & 1.0000 & 1{,}399{,}040 & 1{,}399{,}040 & 14.4\% \\
 & 4 & 148{,}856 & 148{,}856 & 1.0000 & 716{,}160 & 716{,}160 & 0.0\% \\
 & 5 & 120{,}170 & 120{,}170 & 1.0000 & 954{,}880 & 954{,}880 & 0.0\% \\
\addlinespace[2pt]
case\_090 & 1 & 842{,}323 & 842{,}323 & 1.0000 & 0 & 0 & 0.0\% \\
 & 2 & 451{,}762 & 451{,}762 & 1.0000 & 2{,}165{,}760 & 2{,}165{,}760 & 36.0\% \\
 & 3 & 333{,}250 & 334{,}165 & 0.9973 & 4{,}521{,}600 & 4{,}521{,}600 & 8.9\% \\
 & 4 & 306{,}241 & 272{,}968 & 1.1219 & 6{,}670{,}080 & 7{,}728{,}000 & 43.7\% \\
 & 5 & 276{,}562 & 206{,}001 & 1.3425 & 8{,}593{,}920 & 8{,}593{,}920 & 50.3\% \\
\addlinespace[2pt]
case\_091 & 1 & 9{,}186{,}769 & 9{,}168{,}820 & 1.0020 & 53{,}947{,}008 & 53{,}947{,}008 & 10.5\% \\
 & 2 & 4{,}588{,}580 & 4{,}579{,}875 & 1.0019 & 42{,}713{,}056 & 46{,}943{,}200 & 36.1\% \\
 & 3 & 3{,}069{,}210 & 3{,}051{,}314 & 1.0059 & 46{,}052{,}672 & 46{,}052{,}672 & 37.3\% \\
 & 4 & 2{,}315{,}629 & 2{,}307{,}833 & 1.0034 & 66{,}681{,}504 & 66{,}681{,}504 & 21.4\% \\
 & 5 & 1{,}872{,}971 & 1{,}864{,}221 & 1.0047 & 53{,}916{,}160 & 53{,}916{,}160 & 43.4\% \\
\addlinespace[2pt]
case\_092 & 1 & 4{,}024{,}721 & 4{,}024{,}721 & 1.0000 & 16{,}013{,}312 & 16{,}013{,}312 & 0.0\% \\
 & 2 & 2{,}012{,}858 & 2{,}010{,}714 & 1.0011 & 8{,}012{,}384 & 8{,}012{,}384 & 14.1\% \\
 & 3 & 1{,}352{,}222 & 1{,}349{,}367 & 1.0021 & 7{,}431{,}360 & 7{,}431{,}360 & 21.0\% \\
 & 4 & 1{,}016{,}169 & 995{,}267 & 1.0210 & 6{,}264{,}608 & 6{,}264{,}608 & 42.1\% \\
 & 5 & 797{,}596 & 789{,}206 & 1.0106 & 3{,}950{,}976 & 4{,}311{,}424 & 57.0\% \\
\addlinespace[2pt]
case\_093 & 1 & 74{,}205 & 74{,}205 & 1.0000 & 0 & 0 & 0.0\% \\
 & 2 & 37{,}827 & 37{,}827 & 1.0000 & 8{,}192 & 8{,}192 & 0.0\% \\
 & 3 & 25{,}103 & 25{,}103 & 1.0000 & 16{,}384 & 16{,}384 & 0.0\% \\
 & 4 & 19{,}779 & 19{,}779 & 1.0000 & 24{,}576 & 24{,}576 & 0.0\% \\
 & 5 & 15{,}876 & 15{,}876 & 1.0000 & 32{,}768 & 32{,}768 & 0.0\% \\
\addlinespace[2pt]
case\_094 & 1 & 93{,}286 & 93{,}286 & 1.0000 & 0 & 0 & 0.0\% \\
 & 2 & 46{,}959 & 46{,}959 & 1.0000 & 106{,}272 & 106{,}272 & 9.6\% \\
 & 3 & 38{,}877 & 38{,}577 & 1.0078 & 1{,}143{,}872 & 1{,}470{,}152 & 15.8\% \\
 & 4 & 26{,}469 & 26{,}469 & 1.0000 & 318{,}816 & 318{,}816 & 25.8\% \\
 & 5 & 24{,}751 & 24{,}751 & 1.0000 & 259{,}200 & 259{,}200 & 14.8\% \\
\addlinespace[2pt]
case\_095 & 1 & 2{,}084{,}853 & 2{,}084{,}853 & 1.0000 & 0 & 0 & 0.0\% \\
 & 2 & 1{,}042{,}890 & 1{,}042{,}890 & 1.0000 & 0 & 0 & 0.0\% \\
 & 3 & 695{,}775 & 695{,}775 & 1.0000 & 0 & 0 & 0.0\% \\
 & 4 & 524{,}422 & 524{,}422 & 1.0000 & 0 & 0 & 0.0\% \\
 & 5 & 420{,}852 & 420{,}852 & 1.0000 & 0 & 0 & 0.0\% \\
\addlinespace[2pt]
case\_096 & 1 & 76{,}553 & 76{,}345 & 1.0027 & 168{,}460 & 168{,}460 & 10.5\% \\
 & 2 & 46{,}739 & 46{,}739 & 1.0000 & 140{,}168 & 140{,}168 & 15.8\% \\
 & 3 & 31{,}518 & 31{,}518 & 1.0000 & 143{,}884 & 143{,}884 & 36.8\% \\
 & 4 & 26{,}984 & 26{,}984 & 1.0000 & 78{,}336 & 78{,}336 & 2.0\% \\
 & 5 & 26{,}984 & 26{,}984 & 1.0000 & 78{,}336 & 78{,}336 & 2.0\% \\
\addlinespace[2pt]
case\_097 & 1 & 11{,}491{,}509 & 11{,}491{,}509 & 1.0000 & 10{,}584{,}064 & 10{,}584{,}064 & 0.0\% \\
 & 2 & 5{,}617{,}602 & 5{,}617{,}602 & 1.0000 & 11{,}796{,}480 & 11{,}796{,}480 & 0.0\% \\
 & 3 & 3{,}660{,}459 & 3{,}660{,}338 & 1.0000 & 16{,}031{,}744 & 16{,}031{,}744 & 22.0\% \\
 & 4 & 2{,}626{,}908 & 2{,}626{,}908 & 1.0000 & 14{,}221{,}312 & 14{,}221{,}312 & 4.2\% \\
 & 5 & 2{,}102{,}348 & 2{,}102{,}250 & 1.0000 & 21{,}311{,}488 & 21{,}311{,}488 & 24.5\% \\
\addlinespace[2pt]
case\_098 & 1 & 177{,}430 & 177{,}430 & 1.0000 & 0 & 0 & 0.0\% \\
 & 2 & 89{,}551 & 89{,}551 & 1.0000 & 0 & 0 & 0.0\% \\
 & 3 & 89{,}551 & 89{,}551 & 1.0000 & 0 & 0 & 0.0\% \\
 & 4 & 88{,}241 & 82{,}689 & 1.0671 & 3{,}435{,}400 & 3{,}435{,}400 & 12.4\% \\
 & 5 & 64{,}329 & 61{,}024 & 1.0542 & 2{,}258{,}230 & 2{,}299{,}158 & 21.3\% \\
\addlinespace[2pt]
case\_099 & 1 & 1{,}622{,}999 & 1{,}618{,}902 & 1.0025 & 5{,}484{,}544 & 5{,}484{,}544 & 25.2\% \\
 & 2 & 804{,}729 & 803{,}485 & 1.0015 & 6{,}524{,}928 & 6{,}524{,}928 & 46.8\% \\
 & 3 & 537{,}579 & 536{,}788 & 1.0015 & 6{,}754{,}304 & 6{,}754{,}304 & 39.0\% \\
 & 4 & 405{,}898 & 402{,}775 & 1.0078 & 9{,}334{,}784 & 9{,}334{,}784 & 53.5\% \\
 & 5 & 331{,}823 & 325{,}561 & 1.0192 & 8{,}712{,}192 & 10{,}092{,}544 & 56.5\% \\
\addlinespace[2pt]
case\_100 & 1 & 151{,}223 & 151{,}223 & 1.0000 & 0 & 0 & 0.0\% \\
 & 2 & 73{,}847 & 73{,}847 & 1.0000 & 72{,}192 & 72{,}192 & 14.2\% \\
 & 3 & 57{,}811 & 57{,}811 & 1.0000 & 72{,}192 & 72{,}192 & 14.2\% \\
 & 4 & 57{,}811 & 57{,}811 & 1.0000 & 72{,}192 & 72{,}192 & 14.2\% \\
 & 5 & 57{,}811 & 57{,}811 & 1.0000 & 72{,}192 & 72{,}192 & 14.2\% \\
\addlinespace[2pt]
\end{longtable}
\end{footnotesize}


\section{关键源程序}

附件程序只依赖 Python 3.10 及以上版本的标准库和赛题官方材料包。下面给出三种构造算法（算法~\ref{alg:lpt}--\ref{alg:phase}）和求解入口的完整源码，求解入口中包含候选生成、官方评估、成组合并与补空核（算法~\ref{alg:select} 的前半部分）。参数扫描、局部搜索、交叉复评与统计作图的实验代码随附件提交，此处从略。

\lstinputlisting[language=Python, title={\small 分量装箱与拓扑区间切分（nmcs/algos.py）}]{code/algos.py}
\lstinputlisting[language=Python, title={\small 层带 $\times$ 核划分（nmcs/algos2.py）}]{code/algos2.py}
\lstinputlisting[language=Python, title={\small 通信感知表调度与相位构造（nmcs/algos4.py）}]{code/algos4.py}
\lstinputlisting[language=Python, title={\small 求解入口（solve.py）}]{code/solve.py}

\section{人工智能工具使用说明}

本说明依据《“华为杯”第二十三届中国研究生数学建模竞赛人工智能工具及输出使用规定》整理。本文的问题理解、建模思路、技术路线以及各项方案取舍由参赛队讨论确定。人工智能工具在队员给定思路和要求后，承担编程实现、批量计算与数据整理、论文语言润色与排版等辅助工作；全部数值均由赛题官方评估程序在固定配置下给出，经队员核验后采用。

\noindent\textbf{一、所用人工智能工具}
\begin{center}\small
\resizebox{\textwidth}{!}{%
\begin{tabular}{cccccc}
\toprule
序号 & 工具名称 & 版本/型号 & 开发机构 & 版本颁布日期 & 使用日期 \\
\midrule
1 & Claude Code & Claude Opus 5.5 & Anthropic & 2026-09-22 & 2026-09-23 至 09-25 \\
2 & OpenAI Codex & GPT-6-Astra、GPT-5.6-Sol & OpenAI & 2026-09-04、2026-07-09 & 2026-09-23 \\
\bottomrule
\end{tabular}}
\end{center}

\noindent\textbf{二、使用环节与目的}
\begin{center}\small
\begin{tabular}{c p{2.6cm} p{9.2cm} c}
\toprule
序号 & 使用环节 & 内容与目的 & 所用工具 \\
\midrule
1 & 编程实现 & 按队员确定的三种构造算法和选优、后处理流程编写与调试代码，整理附件程序 \texttt{solve.py}、\texttt{reproduce.py} & Claude Code \\
2 & 批量计算与数据整理 & 调用官方评估程序完成参数扫描、后处理与终版复评；汇总统计、计算性能参照值与缓存诊断指标，生成数据图 & Claude Code \\
3 & 方案评审与代码检查 & 对问题三的实验方案提出评审意见（成组合并后处理）并做小样本验证；检查实验收尾脚本 & OpenAI Codex \\
4 & 论文润色与排版 & 提供初稿供队员参考，语言润色，LaTeX 与 Word 版式整理和格式检查 & Claude Code \\
\bottomrule
\end{tabular}
\end{center}

\noindent\textbf{三、使用方式}\quad 队员先确定每一步要做什么，再向工具说明具体要求：编程环节给出算法思路、输入输出格式和合法性要求；计算与数据整理环节给出参数范围和统计口径；评审环节提供实验记录，请工具指出方案和代码中可能的问题；写作环节在队员确定的结构和结论之上生成初稿、润色语言和整理版式。建模思路、模型形式、方案取舍和结论不交由工具决定。

\noindent\textbf{四、采纳、人工修改和核验情况}
\begin{center}\small
\begin{tabular}{c c p{11.6cm}}
\toprule
序号 & 所用工具 & 采纳、人工修改和核验情况 \\
\midrule
1 & Claude Code & 队员审阅代码逻辑与论文算法是否一致；三种构造算法的合法性由命题证明，队员逐条核对推导 \\
2 & Claude Code & 所有数值由官方评估程序在固定配置下复评，1{,}400 个主结果与记录逐位一致；工具不产生未经评估程序验证的数字 \\
3 & OpenAI Codex & 成组合并建议经小样本验证、全量复评确认有效后采纳；检查发现的脚本缺陷修复后重新测试 \\
4 & Claude Code & 模型、公式与证明由队员推导核对，正文由队员在初稿基础上改写定稿；参考文献均为正式发表的文献，未引用工具输出 \\
\bottomrule
\end{tabular}
\end{center}

全部程序文件开头均按规定注明了人工智能工具的使用情况。
